% =====================================================================================================
% Online Appendix D: Additional results (modules m1, m2, m3, m4, m5, m7, m8).
% Writer: conclusion writer, 2026-09-24. Tables in paper/tables/appD_*.tex (four generated by
% code/paper/make_appendix_cd_tables.py from output/tables/*.csv; the rest condensed from output/tables/*.tex).
% Numbers and sources: paper/notes/writer_conclusion.md
% =====================================================================================================

\section{Additional Results}\label{app:additional}
\setlength{\emergencystretch}{2em}% allow looser lines in the online appendix

This appendix collects results that support the main text but are not needed to follow it. Sections~\ref{app:add-chinchilla}
through \ref{app:add-flexible} follow the order of the paper: the canonical Chinchilla data, the technology across
sweeps, revealed inference demand, observational production functions, algorithmic progress, bias formulas for input
measurement and flexible inputs, and hyperparameters as flexible inputs. Section~\ref{app:add-claims} reports the verification register for the formal results.

\subsection{The Canonical Chinchilla Data}\label{app:add-chinchilla}

\emph{Exclusion rules.}---Table~\ref{tab:app-selection} shows how the Chinchilla estimates move as the $k$
highest-loss runs are dropped. The data exponent $\beta$ falls almost monotonically, from 0.453 with all 245 runs to
0.347 at $k=15$, while $\alpha$ stays between 0.345 and 0.349. The five runs that \citet{besiroglu2024chinchilla} drop
are the most data-starved runs of the $10^{19}$ IsoFLOP, with log residuals 9 to 45 times the residual standard
deviation. Their influence is leverage, not truncation bias: a truncated-Gaussian likelihood that conditions on the
exclusion rule reproduces the untruncated Gaussian estimates ($\beta=0.406$ in both), and dropping runs by their
data-to-parameter ratio instead of their loss gives the same answer. The revealed-preference verdict at
Chinchilla-70B depends on excluding these outliers but not on how: every rule that removes them yields $w$ between
0.95 and 1.06, whereas keeping them ($k\le2$) gives 1.24 or more and dropping the whole $10^{19}$ budget gives 0.78.

\emph{Duality.}---Table~\ref{tab:app-duality} reports the objects each of Chinchilla's three approaches estimates
and tests the cross-equation restrictions that link them (Proposition~\ref{prop:duality}). Tests must be
design-consistent: comparing the IsoFLOP argmins with the analytic Approach~3 path rejects even our own refit under
the stratified bootstrap, because parabola fits on a finite grid are biased \citep{czech2026problems}. Against the
design-consistent null, our refit and the \citet{besiroglu2024chinchilla} parameters are consistent with the argmins
($p\ge0.16$ in every classical or bootstrap-calibrated variant), whereas the published Approach~3 parameters of \citet{hoffmann2022training} are
marginally inconsistent ($p=0.040$ in the classical test and 0.060 with bootstrap calibration), mainly because their
implied compute-optimal model size at the center of the design is about 12 percent below the argmins. The frontier
restrictions are rejected for every technology under the stratified bootstrap (calibrated $p$ between 0.003 and 0.037)
but not under the wild bootstrap ($p\ge0.15$): the Approach~1 frontier through the minima lies 0.005 to 0.008 nats
below every Approach~3 frontier, a small lack of fit at the IsoFLOP minima.

\emph{Lab-own technologies.}---Table~\ref{tab:app-labs} applies the same comparison to the IsoFLOP experiments that
labs have published. Our Approach~2 fit reproduces Meta's Llama~3 law exactly, but the primal fitted to the same runs
implies a steeper path ($a=0.527$ against 0.463) in which $M^*$ falls rather than rises with compute. The difference
is significant at about the 5 percent level once the primal's own sampling error is included ($t=-2.1$) and not under
the wild bootstrap; the loss units of the digitized figure are unknown. For the Marin ladders the design-consistent
tests do not reject ($p$ between 0.33 and 0.54); about three-quarters of the naive discrepancy is finite-grid parabola
bias and the rest is FLOP accounting, because Marin's nominal budgets differ from $6ND$. The Marin exponents
($\alpha$ between 0.64 and 0.69) are far from Chinchilla's: exponents are specific to the recipe.

\begin{table}[tp]
\centering
\caption{Chinchilla Data: Dropping High-Loss Runs and Other Exclusion Rules}
\label{tab:app-selection}
\footnotesize
\setlength{\tabcolsep}{2pt}
\begin{tabular}{@{}lcccccc@{}}
\toprule
Sample & $n$ & $\alpha$ & $\beta$ & $a$ & $M^*(10^{21})$ & $w$ (70B) \\
\midrule
\multicolumn{7}{@{}l}{\textit{Panel A. Drop the $k$ highest-loss runs (Huber)}} \\
$k=0$ & 245 & 0.349 (0.018) & 0.453 (0.054) & 0.565 (0.032) & 15.5 (3.7) & 1.54 (0.50) \\
$k=1$ & 244 & 0.345 (0.016) & 0.412 (0.041) & 0.544 (0.027) & 17.6 (3.6) & 1.30 (0.31) \\
$k=2$ & 243 & 0.346 (0.016) & 0.403 (0.035) & 0.538 (0.025) & 18.4 (3.5) & 1.24 (0.24) \\
$k=3$ & 242 & 0.347 (0.015) & 0.382 (0.028) & 0.524 (0.022) & 20.0 (3.3) & 1.12 (0.21) \\
$k=4$ & 241 & 0.347 (0.016) & 0.371 (0.024) & 0.516 (0.022) & 21.0 (3.5) & 1.06 (0.19) \\
$k=5$ (Besiroglu) & 240 & 0.347 (0.015) & 0.367 (0.021) & 0.514 (0.020) & 21.4 (3.5) & 1.04 (0.17) \\
$k=6$ & 239 & 0.347 (0.013) & 0.363 (0.018) & 0.511 (0.018) & 21.9 (3.2) & 1.02 (0.13) \\
$k=8$ & 237 & 0.346 (0.013) & 0.358 (0.019) & 0.509 (0.019) & 22.2 (3.5) & 1.00 (0.13) \\
$k=10$ & 235 & 0.346 (0.015) & 0.358 (0.018) & 0.509 (0.018) & 22.2 (3.5) & 1.00 (0.14) \\
$k=15$ & 230 & 0.347 (0.016) & 0.347 (0.020) & 0.500 (0.021) & 23.3 (4.0) & 0.95 (0.14) \\
\addlinespace
\multicolumn{7}{@{}l}{\textit{Panel B. Alternative exclusion rules (Huber)}} \\
Drop 6 lowest-loss runs & 239 & 0.359 (0.014) & 0.479 (0.057) & 0.572 (0.032) & 14.6 (3.6) & 1.67 (0.54) \\
Drop 5 largest residuals & 240 & 0.350 (0.015) & 0.373 (0.024) & 0.516 (0.020) & 21.1 (3.4) & 1.05 (0.18) \\
Drop $D/N<0.4$ & 241 & 0.347 (0.016) & 0.371 (0.024) & 0.516 (0.021) & 21.0 (3.5) & 1.06 (0.18) \\
Drop $D/N<1$ & 234 & 0.343 (0.015) & 0.342 (0.019) & 0.500 (0.021) & 22.7 (3.9) & 0.96 (0.14) \\
Drop $D/N<2$ & 224 & 0.331 (0.013) & 0.347 (0.023) & 0.512 (0.022) & 20.2 (3.4) & 1.05 (0.15) \\
Drop the $10^{19}$ budget & 211 & 0.357 (0.014) & 0.313 (0.022) & 0.467 (0.024) & 27.4 (3.9) & 0.78 (0.12) \\
IsoFLOP-profile runs only & 137 & 0.383 (0.030) & 0.490 (0.055) & 0.561 (0.030) & 15.8 (3.3) & 1.56 (0.58) \\
IsoFLOP-profile runs, $k=5$ & 132 & 0.363 (0.026) & 0.402 (0.028) & 0.526 (0.022) & 18.8 (3.3) & 1.16 (0.19) \\
\addlinespace
\multicolumn{7}{@{}l}{\textit{Panel C. Gaussian likelihood and truncated regression}} \\
Gaussian, all runs & 245 & 0.372 & 0.612 & 0.622 & 9.6 & 3.08 \\
Gaussian, $k=5$ & 240 & 0.360 & 0.406 & 0.530 & 20.8 & 1.14 \\
Truncated Gaussian, $k=5$ & 240 & 0.361 & 0.406 & 0.530 & 20.8 & 1.14 \\
\bottomrule
\end{tabular}
\par\smallskip
\begin{tablenotes}Huber ($\delta=10^{-3}$) estimates on the 245 digitized Chinchilla runs, with pairs-bootstrap standard errors (200 draws) in parentheses; Online Appendix Table~\ref{tab:app-chin-inference} reports design-conditional alternatives for the $k=5$ sample. $w$ is the wedge $\varepsilon_N/\varepsilon_D$ at Chinchilla-70B ($N=70$B, $D=1.4$T), which equals 1 if that model is training-compute-optimal under the estimated technology. Panel A: $k=5$ is the sample of \citet{besiroglu2024chinchilla}; the five dropped runs are the most data-starved runs of the $10^{19}$ IsoFLOP budget ($D/N=0.04$--0.40), with log residuals 9 to 45 times the residual standard deviation. Panel B: selection on the regressor ($D/N$ or budget) instead of the outcome; the IsoFLOP-profile runs are the 137 runs within 0.045 of a nominal budget in $\log_{10}C$. Across the rules that remove the five outliers, $w$(70B) lies between 0.78 and 1.16. Panel C: the truncated-Gaussian likelihood \citep{hausman1977social} conditions on $L$ below the fifth-highest loss (3.447). A paired bootstrap of the $k=0\to5$ change (400 draws, rule applied within each draw) gives $\Delta\beta=0.086$ (s.e.\ 0.052) and $\Delta a=0.051$ (s.e.\ 0.029).\end{tablenotes}
\begin{tablenotes}[Source]Epoch AI digitization of \citet{hoffmann2022training}, Figure 4; authors' estimates.\end{tablenotes}
\end{table}

\begin{table}[tp]
\centering
\caption{Chinchilla Data: Cost Function, Conditional Factor Demand, and Primal Estimates}
\label{tab:app-duality}
\footnotesize
\setlength{\tabcolsep}{3.5pt}
\begin{tabular}{@{}lccccc@{}}
\toprule
\multicolumn{6}{@{}l}{\textit{Panel A. Dual objects implied by each approach}} \\[1pt]
 & IsoFLOP & Minima, & \multicolumn{2}{c}{Approach 3 (primal)} & \\
\cmidrule(lr){4-5}
 & minima & $E$ fixed & Analytic & Design & \\
\midrule
Path slope $a$ & 0.498 (0.020) & 0.498 & 0.514 (0.018) & 0.515 (0.020) & \\
Path level $\ln N^*(\bar C)$ & 20.589 (0.027) & 20.589 & 20.528 (0.036) & 20.560 (0.045) & \\
Frontier elasticity $\gamma$ & 0.168 (0.008) & 0.181 & 0.179 (0.005) & 0.177 (0.006) & \\
Frontier level $L^*(\bar C)$ & 2.557 (0.002) & 2.554 & 2.562 (0.001) & 2.563 (0.002) & \\
\midrule
\multicolumn{6}{@{}l}{\textit{Panel B. Tests of the cross-equation restrictions}} \\[1pt]
 & \multicolumn{2}{c}{Path $(a,\ln N^*)$} & \multicolumn{2}{c}{Frontier $(\gamma,L^*)$} & Classical \\
\cmidrule(lr){2-3}\cmidrule(lr){4-5}
Null technology & $\chi^2_2$ ($p$) & Calibrated $p$ & $\chi^2_2$ ($p$) & Calibrated $p$ & $F_{2,7}$: $p$ \\
\midrule
Our refit, analytic & 6.8 (0.033) & 0.053 / 0.357 & 16.7 ($<$0.001) & 0.007 / 0.489 & -- \\
Our refit, design & 2.9 (0.238) & 0.233 / 0.698 & 20.7 ($<$0.001) & 0.007 / 0.414 & 0.653 \\
Hoffmann, analytic & 53.5 ($<$0.001) & 0.010 / 0.055 & 7.7 (0.022) & 0.037 / 0.848 & -- \\
Hoffmann, design & 15.6 ($<$0.001) & 0.060 / 0.209 & 14.4 ($<$0.001) & 0.017 / 0.696 & 0.040 \\
Besiroglu, analytic & 13.5 (0.001) & 0.027 / 0.406 & 32.0 ($<$0.001) & 0.003 / 0.195 & -- \\
Besiroglu, design & 4.3 (0.118) & 0.176 / 0.746 & 42.5 ($<$0.001) & 0.003 / 0.150 & 0.664 \\
\bottomrule
\end{tabular}
\par\smallskip
\begin{tablenotes}The 240-run Chinchilla sample. IsoFLOP minima: a parabola in $\ln N$ fitted to the 132
IsoFLOP-profile runs of each of the nine budgets, with the argmins regressed on $\ln(C/6)$ (conditional factor demand,
Hoffmann et al.'s Approach 2) and $L^*_b=E+K(C_b/6)^{-\gamma}$ fitted to the nine minima (the IsoFLOP-minima frontier,
an inverse cost function), with $E$ free or fixed at the primal value. This frontier is not Hoffmann et al.'s
Approach 1, which uses envelopes of training curves that the digitization does not contain and which also reports
allocation exponents. Approach 3 is the Huber primal on all 240 runs. Hoffmann and Besiroglu denote the published
parameters of \citet{hoffmann2022training} and \citet{besiroglu2024chinchilla}. ``Analytic'' uses the closed forms
$a=\beta/(\alpha+\beta)$, $\gamma=\alpha\beta/(\alpha+\beta)$, $N^*=G(C/6)^a$ and $L^*=E+K(C/6)^{-\gamma}$; ``design''
(design-consistent) applies the IsoFLOP-minima procedure to the primal's fitted values at the observed design, which
removes the finite-grid bias of parabola fits \citep{czech2026problems}. $\bar C$ is the geometric mean of the budgets.
Standard errors in parentheses: within-budget stratified pairs bootstrap (300 draws). Panel B: Wald statistics with an
outlier-robust bootstrap covariance and $\chi^2_2$ $p$-values; calibrated $p$-values take critical values from the
bootstrap distribution of the same statistic (stratified pairs / fixed-design wild bootstrap); the classical test
regresses the nine argmin gaps on $[1,\ln C_b]$ and treats the null technology as known. For published parameters only
the IsoFLOP-minima side is random. With all 245 runs, the design-consistent path test of our refit gives a calibrated
$p$ of 0.054 (stratified) and a classical $p$ of 0.041.\end{tablenotes}
\begin{tablenotes}[Source]Epoch AI digitization of \citet{hoffmann2022training}; authors' estimates.\end{tablenotes}
\end{table}

\begin{table}[tp]
\centering
\caption{Lab-Own Technologies from Published IsoFLOP Experiments}
\label{tab:app-labs}
\footnotesize
\setlength{\tabcolsep}{3pt}
\begin{tabular}{@{}lccccc@{}}
\toprule
 & Llama 3 & \multicolumn{3}{c}{Marin (2026), by corpus} & (Mis)Fitting \\
\cmidrule(lr){3-5}
 & (Meta) & Comma & DCLM & Nemotron-CC & FineWeb/C4 \\
\midrule
\multicolumn{6}{@{}l}{\textit{Panel A. Primal (Approach 3, Huber) on the IsoFLOP runs}} \\
Runs (budgets) & 133 (10) & 85 (7) & 85 (7) & 88 (8) & 176 (26) \\
$\alpha$ & 0.284 (0.050) & 0.666 (0.085) & 0.686 (0.061) & 0.639 (0.079) & 0.610 (0.051) \\
$\beta$ & 0.316 (0.024) & 0.350 (0.038) & 0.405 (0.034) & 0.400 (0.040) & 0.705 (0.040) \\
$a$ & 0.527 (0.032) & 0.345 (0.026) & 0.371 (0.018) & 0.385 (0.020) & 0.536 (0.019) \\
$\sigma^*$ & 0.769 (0.018) & 0.663 (0.022) & 0.647 (0.017) & 0.658 (0.022) & 0.603 (0.014) \\
$M^*(10^{21})$ & 13.9 (2.7) & 26.6 (5.3) & 17.9 (2.7) & 16.7 (2.7) & 13.8 (3.1) \\
$M^*(3.8\times10^{25})$ & 7.9 & 707 & 270 & 189 & 6.4 \\
\addlinespace
\multicolumn{6}{@{}l}{\textit{Panel B. IsoFLOP argmins and IsoFLOP-minima frontier ($\alpha$, $\beta$ via $\kappa=1$)}} \\
$a$ & 0.463 (0.006) & 0.405 (0.032) & 0.424 (0.030) & 0.422 (0.027) & 0.517 (0.041) \\
$\gamma$ & 0.136 (0.003) & 0.051 (0.011) & 0.076 (0.023) & 0.075 (0.020) & 0.165 (0.078) \\
$M^*(3.8\times10^{25})$ & 40.9 & 116 & 56.7 & 61.8 & 10.2 \\
\addlinespace
\multicolumn{6}{@{}l}{\textit{Panel C. Tests of the IsoFLOP argmins against the primal: statistic ($p$-value)}} \\
Path, design-consistent $\chi^2_2$ & 16.9 ($<$0.001) & 2.2 (0.332) & 1.5 (0.466) & 1.3 (0.535) & 4.5 (0.108) \\
Path, calibrated $p$ & 0.015 & 0.517 & 0.507 & 0.602 & 0.149 \\
Slope, $t$ with primal s.e. & $-$2.11 (0.035) & 0.17 (0.867) & 0.13 (0.895) & $-$0.13 (0.896) & 0.88 (0.378) \\
System LR, small-sample $p$ & 0.050 & 0.246 & 0.391 & 0.351 & 0.135 \\
\bottomrule
\end{tabular}
\par\smallskip
\begin{tablenotes}IsoFLOP runs compiled by \citet{czech2026problems}: Meta's Llama 3 IsoFLOPs digitized from
\citet{grattafiori2024llama} ($N=C/(6D)$; loss units not stated), the Marin Llama-2-architecture ladders on three
corpora \citep{marin2026ladders} ($6ND$ differs from the nominal budgets by $-7$ to $+35$ percent), and the
checkpoint-interpolated IsoFLOPs of \citet{li2025misfitting}. Loss levels are not comparable across columns;
exponents are. Panel A: cluster-bootstrap standard errors over budgets (200 draws). Panel B: robust within-budget
stratified-bootstrap standard errors; the classical standard error of the Llama 3 path slope is 0.018. Our
argmin fit reproduces Meta's published law, $D^*=0.2994\,C^{0.5368}$ against $0.299\,C^{0.537}$, although two of Meta's ten budgets ($3\times10^{21}$ and $10^{22}$ FLOP) do not bracket the minimum with two runs on each side. $M^*$ at
$3.8\times10^{25}$ FLOP extrapolates 3.5 (Llama 3) to 5 (Marin) orders of magnitude beyond the largest budget and is
shown only to illustrate how far the primal and dual extrapolations diverge. Panel C: the calibrated $p$-value
takes critical values from the bootstrap distribution of the path statistic; the slope $t$ adds the primal's
bootstrap standard error to the classical standard error of the argmin slope; the small-sample $p$-value
calibrates the nested system likelihood ratio (level and slope shift of the argmin equation) with an $F$
form.\end{tablenotes}
\begin{tablenotes}[Source]\citet{czech2026problems} and \citet{czech2026llama3isoflop}; authors'
estimates.\end{tablenotes}
\end{table}


\clearpage

\subsection{The Technology across Sweeps}\label{app:add-tech}

\emph{Robustness.}---Table~\ref{tab:app-robust} reports the main estimates of Table~\ref{tab:technology} under
alternative parameter-count conventions, outputs, samples, and estimators. Conventions matter most for the allocation
exponent: counting Farseer's two untied embedding matrices moves $a$ from 0.526 to 0.411 (and $\sigma^*$ from 0.772 to
0.724), and dropping the two single-epoch runs in which \citet{muennighoff2023scaling} trained 1.1 and 2.8 billion
parameters on 100 million tokens moves $a$ from 0.535 to 0.360 and $M^*(10^{21})$ from 60 to 139. The output measure
also matters: with task bits per byte instead of C4 cross-entropy, the OLMo ladder's $\sigma^*$ falls from 0.825 to
0.763. Because $\sigma$ is ordinal, such differences reflect sampling noise and misspecification of the functional
form rather than the change of output units itself.

\emph{Neutrality of data quality.}---Table~\ref{tab:app-neutrality} nests the hypotheses about how better data enter.
Across the three web corpora of \citet{gadre2024language}, a Hicks-neutral scaling of reducible loss plus
corpus-specific asymptotes is not rejected ($p=0.90$) and explains 96.8 percent of the cross-corpus variation, but the
test has little power: the data- and parameter-augmenting alternatives are not rejected either. Across the 25
DataDecide recipes \citep{magnusson2025datadecide} every restriction is rejected, yet asymptote shifts alone explain
98.1 percent of the variation and a neutral scale brings this to 99.7 percent. The remaining factor bias is small in fit
but not in allocation: the recipe tilts $\ln(A_r/B_r)$ span 0.22 to 0.26, which implies compute-optimal ratios that
differ 2.9- to 3.4-fold across recipes and revealed wedges that differ by a factor of 1.25 to 1.29 at the same inputs.
The data-augmenting tilt belongs to the heavily filtered DCLM variants, which have the highest C4 asymptotes. Because
a data-augmenting tilt lowers $M^*$ by construction when exponents are common, and because quality is not measured
independently of the tilt, this pattern is not independent evidence on whether better data are data-augmenting.

\emph{Functional form on the best-designed sweep.}---Table~\ref{tab:app-farseer} compares the Chinchilla form with
richer forms on the Farseer grid of \citet{li2025predictableb}. The Chinchilla form extrapolates poorly in compute:
its out-of-sample error on the top decile of compute is 8 to 10 times that of Farseer's own nine-parameter form.
Figure~\ref{fig:app-farseer-sigma} shows the local elasticity of substitution estimated without any functional form.
Its median over the grid, 0.690, agrees with Farseer's form (0.706) and with the generalized family with a free outer
exponent (0.710), and lies below the Chinchilla-implied 0.77. Local $\sigma$ tends to fall with $D/N$ for the smallest
models, a pattern suggestive of stronger complementarity where over-trained models sit; it is not monotone at every
size, and the pointwise bands ignore smoothing bias.

\begin{table}[tp]
\centering
\caption{Robustness of the Technology Estimates across Sweeps}
\label{tab:app-robust}
\footnotesize
\setlength{\tabcolsep}{4pt}
\begin{tabular}{@{}lcccc@{}}
\toprule
 & \multicolumn{3}{c}{Huber} & Gaussian NLS \\
\cmidrule(lr){2-4}\cmidrule(l){5-5}
Sweep and specification & $\sigma^*$ & $a$ & $M^*(10^{21})$ & $\sigma^*$ \\
\midrule
Farseer: baseline & 0.772 (0.010) & 0.526 (0.022) & 24.8 & 0.769 (0.006) \\
\quad $N$ including embeddings & 0.724 (0.011) & 0.411 (0.017) & 19.8 & 0.721 (0.007) \\
\quad Training loss (nats per token) & 0.785 (0.012) & 0.535 (0.026) & 18.1 & 0.786 (0.007) \\
\quad Clusters = model size & 0.772 (0.009) & 0.526 (0.014) & 24.8 & 0.769 (0.005) \\
\addlinespace
Gadre et al., C4: baseline & 0.795 (0.020) & 0.509 (0.072) & 6.9 & 0.803 (0.020) \\
\quad In-distribution evaluation & 0.799 (0.019) & 0.487 (0.078) & 9.8 & 0.806 (0.017) \\
\quad Non-embedding $N$ & 0.806 (0.016) & 0.534 (0.069) & 11.0 & 0.816 (0.018) \\
Gadre et al., RedPajama: baseline & 0.795 (0.013) & 0.529 (0.045) & 5.8 & 0.809 (0.014) \\
\quad In-distribution evaluation & 0.753 (0.018) & 0.559 (0.083) & 8.0 & 0.772 (0.018) \\
\quad Non-embedding $N$ & 0.808 (0.015) & 0.586 (0.042) & 4.8 & 0.821 (0.014) \\
Gadre et al., RefinedWeb: baseline & 0.796 (0.015) & 0.566 (0.055) & 3.4 & 0.806 (0.016) \\
\quad In-distribution evaluation & 0.805 (0.018) & 0.564 (0.055) & 2.8 & 0.817 (0.015) \\
\quad Non-embedding $N$ & 0.806 (0.014) & 0.606 (0.059) & 3.8 & 0.818 (0.014) \\
\addlinespace
OLMo ladder: baseline & 0.825 (0.034) & 0.370 (0.177) & 48.3 & 0.833 (0.024) \\
\quad Task bits per byte as output & 0.763 (0.030) & 0.353 (0.116) & 85.4 & 0.770 (0.033) \\
\addlinespace
Muennighoff et al.: single epoch & 0.735 (0.031) & 0.535 (0.126) & 59.6 & 0.723 (0.031) \\
\quad Single epoch, $D/N\geq0.4$ & 0.726 (0.016) & 0.360 (0.086) & 139 & 0.711 (0.018) \\
\quad Up to 4 epochs & 0.732 (0.014) & 0.522 (0.067) & 71.6 & 0.719 (0.014) \\
\bottomrule
\end{tabular}
\par\smallskip
\begin{tablenotes}Chinchilla functional form ($\kappa=1$); bootstrap standard errors in parentheses (400 draws for
baselines, 200 for variants). Farseer \citep{li2025predictableb}: $N$ including embeddings adds the two untied
65,536-row embedding matrices to the non-embedding count; training loss replaces English bits per character.
\citet{gadre2024language}: in-distribution evaluation uses the training corpus's Paloma subset instead of C4. OLMo
ladder \citep{bhagia2024establishing}: task bits per byte instead of C4 cross-entropy. \citet{muennighoff2023scaling}:
$D/N\geq0.4$ drops the two single-epoch runs trained on 100 million tokens with 1.1 and 2.8 billion parameters, the
exclusion rule that \citet{besiroglu2024chinchilla} apply to the Chinchilla data; up to 4 epochs treats repeated
passes as fresh tokens.\end{tablenotes}
\begin{tablenotes}[Source]Run-level data released with each sweep; authors' estimates.\end{tablenotes}
\end{table}

\begin{table}[tp]
\centering
\caption{Is Better Data Hicks-Neutral? Nested Tests across Corpora and Recipes}
\label{tab:app-neutrality}
\footnotesize
\setlength{\tabcolsep}{4pt}
\begin{tabular}{@{}lccccc@{}}
\toprule
Model for corpus or recipe $r$ & Parameters & Share explained (\%) & LR & df & $p$ (wild) \\
\midrule
\multicolumn{6}{@{}l}{\textit{Panel A. \citet{gadre2024language}: three web corpora (104 runs, 37 clusters)}} \\
(0) One technology & 5 & 0.0 & 30.7 & 6 & $\leq$0.002 \\
(iv) Asymptote shift $E_r$ only & 7 & 89.9 & 2.6 & 4 & 0.004 \\
(i) Hicks-neutral, common $E$ & 7 & 72.7 & 8.6 & 4 & $\leq$0.002 \\
(ii) Data-augmenting, common $E$ & 7 & 64.1 & 11.4 & 4 & $\leq$0.002 \\
(iii) Parameter-augmenting, common $E$ & 7 & 75.0 & 7.8 & 4 & $\leq$0.002 \\
(i$'$) Hicks-neutral + $E_r$ & 9 & 96.8 & 0.0 & 2 & 0.896 \\
(ii$'$) Data-augmenting + $E_r$ & 9 & 95.4 & 0.5 & 2 & 0.130 \\
(iii$'$) Parameter-augmenting + $E_r$ & 9 & 95.9 & 0.4 & 2 & 0.150 \\
Common exponents ($A_r$, $B_r$, $E_r$ free) & 11 & 96.9 & 1.2 & 4 & 0.090 \\
(v) Exponents vary & 15 & 100.0 & & & \\
\addlinespace
\multicolumn{6}{@{}l}{\textit{Panel B. DataDecide: 25 recipes (21,888 checkpoints, 44 clusters)}} \\
(0) One technology & 5 & 0.0 & 58,185 & 72 & $\leq$0.050 \\
(iv) Asymptote shift $E_r$ only & 29 & 98.1 & 4,465 & 48 & $\leq$0.010 \\
(i) Hicks-neutral, common $E$ & 29 & 90.4 & 17,732 & 48 & $\leq$0.050 \\
(ii) Data-augmenting, common $E$ & 29 & 92.6 & 14,600 & 48 & $\leq$0.050 \\
(iii) Parameter-augmenting, common $E$ & 29 & 89.4 & 18,964 & 48 & $\leq$0.050 \\
(i$'$) Hicks-neutral + $E_r$ & 53 & 99.7 & 145 & 24 & $\leq$0.010 \\
(ii$'$) Data-augmenting + $E_r$ & 53 & 99.7 & 69 & 24 & $\leq$0.010 \\
(iii$'$) Parameter-augmenting + $E_r$ & 53 & 99.5 & 738 & 24 & $\leq$0.010 \\
Common exponents ($A_r$, $B_r$, $E_r$ free) & 77 & 99.8 & 705 & 48 & $\leq$0.010 \\
(v) Exponents vary & 125 & 100.0 & & & \\
\bottomrule
\end{tabular}
\par\smallskip
\begin{tablenotes}All models are evaluated on the same C4 validation set. Models nest in
$L_r=E_r+A_rN^{-\alpha_r}+B_rD^{-\beta_r}$: (i) Hicks-neutral in reducible loss, $A_r=e^{-\omega_r}A$ and
$B_r=e^{-\omega_r}B$; (ii) only $B_r$ varies; (iii) only $A_r$ varies; (iv) only $E_r$ varies; primes add
recipe-specific $E_r$. Share explained $=(\mathrm{SSR}_0-\mathrm{SSR}_m)/(\mathrm{SSR}_0-\mathrm{SSR}_v)$, where model
(0) is one common technology and (v) separate fits. LR $=n\ln(\mathrm{SSR}_m/\mathrm{SSR}_{CE})$ against the
common-exponent model (for that row, against (v)), from Gaussian NLS on log loss. $p$-values from a wild cluster
restricted bootstrap with Rademacher weights (clusters: size $\times$ token-multiplier cells for Gadre et al., 499
draws; size $\times$ seed cells for DataDecide, 99 draws for the key tests and 19 for the common-$E$ models;
$\leq$ marks the smallest attainable value). Wald tests with a cell-level pairs-bootstrap covariance confirm every
DataDecide rejection and, for Gadre et al., the non-rejection of (i$'$) and the rejection of (iv); they disagree
across estimators on (ii$'$). DataDecide uses intermediate checkpoints with $D/N\geq5$, evaluated before the
learning-rate schedule has finished; the comparisons across recipes assume that this schedule effect is common to all
recipes.\end{tablenotes}
\begin{tablenotes}[Source]\citet{gadre2024language}; \citet{magnusson2025datadecide}; authors'
estimates.\end{tablenotes}
\end{table}

\begin{table}[tp]
\centering
\caption{Farseer: The Chinchilla Form against Richer Functional Forms}
\label{tab:app-farseer}
\footnotesize
\setlength{\tabcolsep}{3pt}
\begin{tabular}{@{}lccccc@{}}
\toprule
 & & \multicolumn{2}{c}{In sample} & \multicolumn{2}{c}{Out of sample: RMSE $\times10^3$} \\
\cmidrule(lr){3-4}\cmidrule(l){5-6}
Model & $k$ & RMSE $\times10^3$ & BIC & Largest $N$ & Top $C$ decile \\
\midrule
Chinchilla, Huber & 5 & 8.74 & $-$2,647 & 12.65 & 24.75 \\
Chinchilla, NLS & 5 & 8.50 & $-$2,670 & 13.79 & 21.95 \\
$\kappa$ free, NLS & 6 & 2.47 & $-$3,663 & 3.57 & 5.23 \\
Translog, $E$ free, NLS & 7 & 2.05 & $-$3,806 & 3.32 & 4.69 \\
Farseer Eq.~3, NLS & 9 & 1.76 & $-$3,916 & 2.63 & 2.58 \\
\bottomrule
\end{tabular}
\par\smallskip
\begin{tablenotes}The Farseer grid of \citet{li2025predictableb}: 404 runs, output English bits per character; RMSE of
log loss; $k$ is the number of parameters. The translog is quadratic in $(\ln N,\ln D)$ for $\ln(L-E)$. Farseer's Eq.~3 is
$L=e^{a_3N^{\gamma_3}+b_3}+e^{a_2N^{\gamma_2}+b_2}D^{-e^{a_1N^{\gamma_1}+b_1}}$, fitted here by NLS on log loss
rather than the authors' differential procedure. Out-of-sample fits start only from training-sample values; the
splits hold out the four largest model sizes ($N\geq3.8\times10^9$) or the top decile of $6ND$. The like-for-like model-free comparison is on the local expansion path ($w=1$): its first-derivative estimate of
$\sigma^*$ is 0.708 (wild cluster standard error 0.001; 0.701--0.711 across bandwidths), against 0.710 for the
$\kappa$-free form and 0.772 for the Chinchilla form; a Hessian-based estimate on the same path is 0.703 (0.013) at the
primary bandwidth but moves between 0.66 and 0.73 with the bandwidth. Eq.~3 implies $\sigma^*(C)$ falling from 0.745 at
$10^{18}$ to 0.672 at $10^{22}$ FLOP. Eq.~3 implies a compute-optimal ratio $M^*$ of 32, 25, 28, 43, and 82 at
$10^{19}$ to $10^{23}$ FLOP; the last two are extrapolations.\end{tablenotes}
\begin{tablenotes}[Source]\citet{li2025predictableb}; authors' estimates.\end{tablenotes}
\end{table}


\begin{figure}[tp]
\centering
\includegraphics[width=\textwidth]{m2_farseer_local_sigma}
\caption{Local Elasticity of Substitution on the Farseer Grid}\label{fig:app-farseer-sigma}
\begin{figurenotes}
Local $\sigma$ against tokens per parameter for four model sizes. Solid: kernel-weighted local quadratic in
$(\ln N,\ln D)$ fitted to log loss, with a 95 percent pointwise pairs-bootstrap band at a fixed cross-validated
bandwidth; the band ignores smoothing bias. Dashed: $\sigma$ implied by the Chinchilla fit. Dotted: $\sigma$ implied
by Farseer's Eq.~3. Because $\sigma$ is invariant to monotone transformations of output, the local estimate needs no
estimate of the irreducible loss $E$. Over 42 interior grid points the median local $\sigma$ is 0.690 (bootstrap 95
percent interval $[0.679,0.699]$), with a median of 0.694 and 0.707 at 1.5 and 2 times the cross-validated bandwidth.
Source: \citet{li2025predictableb}; authors' estimates.
\end{figurenotes}
\end{figure}

\clearpage

\subsection{Revealed Inference Demand}\label{app:add-wedge}

\emph{Families.}---Table~\ref{tab:app-family} normalizes each family by its flagship. Smaller siblings are more
over-trained than their flagships in essentially every family: 95 percent of the 105 non-flagship siblings have
$w^{rel}>1$, and the exceptions are the five smaller Cerebras-GPT models, trained at a constant 20 tokens per
parameter. This pattern follows largely from the exponents once $D$ is given. The assumption that flagships are
training-optimal holds before 2024 (median flagship wedge 1.14; the 95 percent interval contains 1 for half of the 18
flagships) but fails afterward (median 3.14; the interval lies above 1 for 92 percent of 26 flagships), so
flagship-normalized demands are lower bounds for recent families. Meta is the exception: its flagships have $w=0.98$
under Meta's own IsoFLOP law. A single power-law demand for model size across families is rejected, but the test is
partly mechanical when $D$ is common within a family; among families whose token counts vary with size the rejection is
marginal ($p=0.047$) and disappears with flagship-normalized demand ($p=0.12$).

\emph{Rival explanations.}---Table~\ref{tab:app-rivals} evaluates alternatives to inference demand. Under-training
($w<1$), the signature of data scarcity or of Kaplan-era beliefs, disappears after 2023. Hardware tiers matter for
model sizes (26 percent of open models released since 2023 sit at 6.5 to 9.5 billion parameters, against 9 percent
under a log-normal benchmark) but only marginally for the wedge ($+0.11$ in logs, s.e.\ 0.06). A non-homothetic
technology can rationalize the token counts of the largest models but not those of the small ones: within Farseer's
support its Eq.~3 implies larger wedges than the Chinchilla form. Recipe-level factor bias of the size measured in
DataDecide cannot bring most models to $w\le1$. Table~\ref{tab:app-aggregate} reports the compute-weighted aggregate
by release year, and Table~\ref{tab:app-homothetic} shows how revealed $T/D$ for five prominent models depends on the
assumed compute-optimal ratio and curvature when $\alpha=\beta$.

\begin{table}[tp]
\centering
\caption{Within-Family Revealed Preference and Over-Identification}
\label{tab:app-family}
\footnotesize
\setlength{\tabcolsep}{3pt}
\begin{tabular}{@{}lcccccc@{}}
\toprule
 & & \multicolumn{4}{c}{Flagship} & Siblings \\
\cmidrule(lr){3-6}
Family (flagship size) & Models & $M_f$ & $\hat w_f$ & Band & Meta law & $w^{rel}_i>1$ \\
\midrule
\multicolumn{7}{@{}l}{\textit{Panel A. Families}} \\
Llama 2 (70B) & 3 & 29 & 1.19 & [0.55, 3.77] & 0.98 & 1.00 \\
Llama 3.1 (405B) & 3 & 38 & 1.37 & [0.42, 5.73] & 0.98 & 1.00 \\
Qwen2.5 (72B) & 7 & 248 & 2.62 & [1.06, 8.01] & 1.69 & 1.00 \\
Qwen3 (32B) & 6 & 1,099 & 4.46 & [1.98, 11.86] & 2.55 & 1.00 \\
Gemma 2$^{\dagger}$ (27B) & 3 & 477 & 3.27 & [1.58, 8.74] & 2.08 & 1.00 \\
Gemma 3$^{\dagger}$ (27B) & 5 & 518 & 3.37 & [1.62, 8.98] & 2.12 & 1.00 \\
OLMo 2$^{\dagger}$ (32B) & 4 & 191 & 2.34 & [0.23, 6.42] & 1.64 & 1.00 \\
SmolLM2$^{\dagger}$ (1.7B) & 3 & 6,428 & 8.05 & [4.01, 14.97] & 4.50 & 1.00 \\
Phi-3 (14B) & 3 & 344 & 2.86 & [1.63, 6.98] & 1.97 & 1.00 \\
Pythia (12B) & 8 & 25 & 1.09 & [0.75, 2.66] & 1.02 & 1.00 \\
Cerebras-GPT$^{\dagger}$ (6.7B) & 6 & 20 & 0.99 & [0.78, 2.24] & 0.98 & 0.00 \\
BLOOM (176B) & 5 & 2.1 & 0.46 & [0.20, 1.78] & 0.48 & 1.00 \\
OPT (66B) & 8 & 4.5 & 0.60 & [0.31, 2.00] & 0.61 & 1.00 \\
\addlinespace
\multicolumn{7}{@{}l}{\textit{Panel B. Demand model $T_i=\kappa_fN_i^{-\eta}$ with family effects}} \\[2pt]
 & Families & Models & $\hat\eta$ (s.e.) & $F$ & $p$ & \\
\cmidrule(lr){2-6}
Absolute $T_i$, all families & 42 & 138 & 0.53 (0.07) & 2.55 & 0.001 & \\
Absolute $T_i$, size-varying $D$ & 18 & 58 & 0.36 (0.12) & 2.15 & 0.047 & \\
Normalized $T_i$, all families & 23 & 80 & 0.58 (0.05) & 2.77 & 0.004 & \\
Normalized $T_i$, size-varying $D$ & 11 & 33 & 0.48 (0.12) & 2.08 & 0.122 & \\
\bottomrule
\end{tabular}
\par\smallskip
\begin{tablenotes}Families of verified open-weight base models trained by one developer on one data release; the
flagship is the largest model by total parameters. $\hat w_f$: flagship wedge under the reference technology (our
Huber refit of the Chinchilla data); band: range of 95 percent bootstrap intervals and point estimates over six
technologies (a sensitivity band, not an identified set); Meta law: $\hat w_f$ under Meta's Llama 3 IsoFLOP technology. Siblings: share of non-flagship models with $w^{rel}_i>1$. Siblings'
wedges relative to the flagship, $w^{rel}_i=(N_i/N_f)^{-\alpha}(D_i/D_f)^{\beta}$, depend only on the exponents if
the flagship is training-optimal and are lower bounds on $w_i$ when $w_f>1$. $^{\dagger}$: the largest and smallest token counts
in the family differ by more than 10 percent. Panel B: family-clustered standard errors in
parentheses; $F$ tests a common $\eta$ across families (classical, models with $T_i>0$ only); normalized demand divides
by the flagship's. With a common $D$ within a
family, $T_i/D$ is a mechanical function of $N_i$, so the rows with size-varying $D$ are the informative
ones.\end{tablenotes}
\begin{tablenotes}[Source]Model cards, technical reports, and Hugging Face configuration files; authors'
calculations.\end{tablenotes}
\end{table}

\begin{table}[tp]
\centering
\caption{Rival Explanations for Over-Training and Stated-Intent Checks}
\label{tab:app-rivals}
\footnotesize
\setlength{\tabcolsep}{3pt}
\begin{tabular}{@{}>{\raggedright\arraybackslash}p{0.23\textwidth}>{\raggedright\arraybackslash}p{0.43\textwidth}>{\raggedright\arraybackslash}p{0.30\textwidth}@{}}
\toprule
Rival explanation & Statistic & Estimate \\
\midrule
Data scarcity or Kaplan-era beliefs & Share of production-scale models with $\hat w<1$, by release year & 2021: 0.47; 2022: 0.41; 2023: 0.18; 2024: 0.00; 2025: 0.02 \\[2pt]
Memory and hardware tiers & Share of 2023--2026 open models with 6.5--9.5B parameters; log-normal benchmark & 0.26 against 0.09 \\
 & $\Delta\ln\hat w$ for tier models, given $\ln C$ and year & 0.11 (0.06); $n=238$ \\[2pt]
Non-homothetic technology (Farseer Eq.~3) & Median $\hat w$, Chinchilla form against Eq.~3: $N$ inside Farseer's range & 3.44 against 4.12 ($n=66$) \\
 & The same, $N$ above Farseer's range & 2.12 against 1.54 ($n=75$) \\[2pt]
Distillation and synthetic data & $\Delta\ln\hat w$ given $\ln C$ and year: distilled; synthetic & 0.02 (0.23); 0.06 (0.18); $n=173$ \\[2pt]
Factor-biased recipes & Share of models with $\hat w$ / lower band above the recipe wedge ratio: 1.29 (DataDecide, main) & 0.88 / 0.69 \\
 & The same, ratio 1.80 (largest robustness variant) & 0.79 / 0.49 \\
\midrule
\multicolumn{3}{@{}l}{\textit{Stated intent}} \\
Chinchilla & Compute-optimal by design \citep{hoffmann2022training} & $\hat w$ 1.04; Meta law 0.89 \\
Cerebras-GPT, 111M--13B & 20 tokens per parameter, stated as compute-optimal \citep{dey2023cerebras} & $\hat w$ 0.92--1.00; Meta law 0.95--1.16 \\
Pythia, 70M--12B & Same 300B tokens at every size by research design \citep{biderman2023pythia} & $\hat w$ 1.09--6.49; Meta law 1.02--4.61 \\
LLaMA 7B, 13B & Best performance at given inference budgets \citep{touvron2023llama} & $\hat w$ 1.65--2.07; Meta law 1.35--1.63 \\
Llama 3 8B, 70B & Small models trained far beyond compute-optimal for inference \citep{grattafiori2024llama} & $\hat w$ 2.48--5.27; Meta law 1.63--3.09 \\
\bottomrule
\end{tabular}
\par\smallskip
\begin{tablenotes}Production-scale models have $6ND\geq10^{21}$ FLOP (335 models, 2019--2026); other rows use the 173
verified open-weight base models. Standard errors, clustered by developer, in parentheses. Memory tier: 6.5 to 9.5
billion total parameters, the range that fits one 16--24 GB accelerator in 16-bit precision; the benchmark is the share
implied by a normal fit to $\log_{10}N$. Farseer's range is non-embedding $N\leq6.4$ billion; above it Eq.~3 is
extrapolated one to two orders of magnitude. Factor bias: the ratio $\exp(\text{tilt range})$ across the 25 DataDecide
recipes bounds how far a recipe's factor bias can move $\hat w$ (Proposition~\ref{prop:wedge}); a data constraint lowers
$w$, a memory constraint raises it. Stated intent: $\hat w$ under the reference technology; Meta law: Meta's Llama 3
IsoFLOP technology.\end{tablenotes}
\begin{tablenotes}[Source]Epoch AI model database \citep{epochai2026data}; model cards and technical reports; authors'
calculations.\end{tablenotes}
\end{table}

\begin{table}[tp]
\centering
\caption{Planned Lifetime Inference Compute of the Open-Weight Ecosystem}
\label{tab:app-aggregate}
\footnotesize
\setlength{\tabcolsep}{3pt}
\begin{tabular}{@{}lcccccc@{}}
\toprule
 & & & & \multicolumn{3}{c}{Inference compute / training compute} \\
\cmidrule(l){5-7}
Year & Models & $\sum 6ND$ & Median $\hat w$ & Estimate [95\% CI] & Band & $T\geq0$ \\
\midrule
2021 & 10 & $1.1\times10^{23}$ & 1.93 & 0.63 [0.38, 1.05] & [0.20, 2.58] & 0.65 \\
2022 & 21 & $1.6\times10^{24}$ & 1.07 & $-$0.40 [$-$0.50, $-$0.23] & [$-$0.68, 0.99] & 0.03 \\
2023 & 81 & $1.4\times10^{25}$ & 2.00 & 0.53 [0.20, 1.15] & [$-$0.18, 3.54] & 0.54 \\
2024 & 95 & $1.2\times10^{26}$ & 3.66 & 1.16 [0.60, 2.27] & [$-$0.04, 5.75] & 1.16 \\
2025 & 58 & $9.2\times10^{25}$ & 5.01 & 3.22 [2.11, 5.67] & [1.05, 9.64] & 3.22 \\
2026 & 4 & $2.0\times10^{25}$ & 3.77 & 2.78 [1.76, 5.10] & [0.65, 9.39] & 2.78 \\
\addlinespace
All years & 271 & $2.5\times10^{26}$ & 2.86 & 2.00 [1.23, 3.69] & [0.44, 7.17] & 2.01 \\
\bottomrule
\end{tabular}
\par\smallskip
\begin{tablenotes}Open-weight models with $6ND\geq10^{21}$ FLOP (one model each in 2019 and 2020, included in the last
row). The multiple $\sum_i2N_iT_i/\sum_i6N_iD_i=\sum_i(w_i-1)C_i/\sum_iC_i$ is the compute-weighted mean of $w-1$:
planned lifetime inference compute as a multiple of training compute. Reference: our Huber refit of the Chinchilla
data, with a 95 percent bootstrap interval. Band: range over the Chinchilla refit, \citet{besiroglu2024chinchilla},
\citet{hoffmann2022training}, Farseer with embeddings, \citet{gadre2024language} RefinedWeb, and Meta's law, including
their bootstrap intervals; across these technologies the all-years point estimate ranges from 0.67 to 3.26.
$T\geq0$ imposed: $T_i$ set to zero when $w_i<1$, since under-training is not negative demand. The multiple is
dominated by the largest models, whose wedges are the most extrapolated, and is not validated against inference
volumes; it illustrates the object rather than measuring it.\end{tablenotes}
\begin{tablenotes}[Source]Epoch AI model database \citep{epochai2026data}; model cards; authors'
calculations.\end{tablenotes}
\end{table}

\begin{table}[tp]
\centering
\caption{Revealed Inference Demand under Homothetic Technologies}
\label{tab:app-homothetic}
\footnotesize
\begin{tabular}{@{}lccccc@{}}
\toprule
 & \multicolumn{5}{c}{Compute-optimal tokens per parameter, $M^*$} \\
\cmidrule(l){2-6}
Model & 16 & 20 & 24.2 & 41 & 192 \\
\midrule
\multicolumn{6}{@{}l}{\textit{Panel: $\rho=0.30$}} \\
Llama 3 8B & 9.5 & 8.7 & 8.1 & 6.4 & 2.9 \\
Llama 3 70B & 3.5 & 3.1 & 2.8 & 1.9 & 0.1 \\
Llama 3.1 405B & 0.9 & 0.6 & 0.4 & $-$0.1 & $-$1.1 \\
Qwen3 8B & 13.2 & 12.1 & 11.3 & 9.2 & 4.7 \\
Gemma 3 27B & 5.5 & 5.0 & 4.5 & 3.4 & 1.0 \\
\addlinespace
\multicolumn{6}{@{}l}{\textit{Panel: $\rho=0.3527$}} \\
Llama 3 8B & 13.1 & 11.9 & 10.9 & 8.5 & 3.7 \\
Llama 3 70B & 4.5 & 3.9 & 3.5 & 2.4 & 0.1 \\
Llama 3.1 405B & 1.1 & 0.8 & 0.5 & $-$0.1 & $-$1.3 \\
Qwen3 8B & 18.7 & 17.1 & 15.8 & 12.6 & 6.1 \\
Gemma 3 27B & 7.2 & 6.5 & 5.8 & 4.3 & 1.3 \\
\addlinespace
\multicolumn{6}{@{}l}{\textit{Panel: $\rho=0.40$}} \\
Llama 3 8B & 17.1 & 15.4 & 14.1 & 10.8 & 4.5 \\
Llama 3 70B & 5.4 & 4.7 & 4.2 & 2.8 & 0.1 \\
Llama 3.1 405B & 1.3 & 0.9 & 0.6 & $-$0.1 & $-$1.4 \\
Qwen3 8B & 25.4 & 22.9 & 21.0 & 16.5 & 7.5 \\
Gemma 3 27B & 9.1 & 8.0 & 7.2 & 5.3 & 1.5 \\
\bottomrule
\end{tabular}
\par\smallskip
\begin{tablenotes}Entries are revealed lifetime inference tokens per training token, $T/D=3[(M/M^*)^{\rho}-1]$, under a homothetic technology with $\alpha=\beta=\rho$, for which $w=(M/M^*)^{\rho}$. $M^*=16$ \citep{porian2024resolving}, 20 (the Chinchilla rule of thumb), 24.2 (our CES fit to the Chinchilla data), 41 (Meta's Llama 3 law at $3.8\times10^{25}$ FLOP), and 192 \citep{hu2024minicpm}; $\rho=0.3527$ is the tied exponent of \citet{muennighoff2023scaling}. Negative entries mean the model is under-trained relative to that technology.\end{tablenotes}
\begin{tablenotes}[Source]Model sizes and token counts from developers' reports and Hugging Face; authors' calculations.\end{tablenotes}
\end{table}


\clearpage

\subsection{Observational Production Functions}\label{app:add-obs}

\emph{Other outputs and a cleaner sample.}---Table~\ref{tab:app-lalonde-robust} repeats the design-matched
comparison of Table~\ref{tab:lalonde} for two other benchmarks and for a sample without imputed token counts,
distilled, synthetic-data, or code models. The OLS bias in the compute elasticity is $-0.098$ (0.127) for ARC-Challenge
and $-0.160$ (0.102) for Winogrande, point estimates that imply 18 to 31 percent attenuation but are not significant; on
the cleaner HellaSwag sample it is $+0.045$ (0.023), significant under clustered normal inference but not under the
wild-cluster bootstrap ($p=0.10$). The parameter--data split is less stable: the sign of the $\theta_N$ bias changes
across outputs and estimators, and the family fixed-effects estimates, which rest on little within-family variation in
$D$, show large negative $\theta_D$ biases for every output.

\emph{Productivity dispersion.}---Table~\ref{tab:app-tfp} and Figure~\ref{fig:app-tfp-families} measure dispersion
in family productivity in compute-equivalent units. The 90/10 ratio across families is 23.7 (90 percent interval
$[9.1,197]$) for HellaSwag, falls to 18.3 without distilled and synthetic-data families and to 7.0 without code
families as well, and is 11.1 within developers, against 41 reported by \citet{mertens2026secret}. The synthetic-data
families Phi and SmolLM rank first and third of 38: when training data are generated by other models, measured
productivity includes the quality of those inputs, a gross-output problem. With returns to scale near one,
manufacturing's 90/10 ratio of 1.92 \citep{syverson2011determines} is also an input-equivalent ratio, so family
dispersion among language models is roughly three to five times larger in logs; expressed in reducible-loss units
(1.41 to 1.76) it looks similar to manufacturing, but that conversion is a choice of output units.

\begin{table}[tp]
\centering
\caption{Design-Matched Benchmark: Other Outputs and a Cleaner Sample}
\label{tab:app-lalonde-robust}
\footnotesize
\setlength{\tabcolsep}{4pt}
\begin{tabular}{@{}lccccc@{}}
\toprule
 & $\theta_N$ & $\theta_D$ & \multicolumn{3}{c}{$\theta_C$ (regression on $\ln C$ only)} \\
\cmidrule(l){4-6}
Estimator & Bias & Bias & Observational & Experimental & Bias \\
\midrule
\multicolumn{6}{@{}l}{\textit{Panel A. ARC-Challenge log odds (57 models)}} \\
OLS & $-$0.153 & $-$0.064 & 0.449 & 0.547 & $-$0.098 \\
 & (0.243) & (0.184) & (0.050) & (0.116) & (0.127) \\
Year FE & $-$0.125 & $-$0.124 & 0.419 & 0.544 & $-$0.124 \\
 & (0.256) & (0.213) & (0.053) & (0.107) & (0.119) \\
Developer FE & 0.097 & $-$0.119 & 0.569 & 0.582 & $-$0.012 \\
 & (0.230) & (0.200) & (0.072) & (0.098) & (0.122) \\
Family FE & 0.235 & $-$0.472 & 0.598 & 0.624 & $-$0.026 \\
 & (0.232) & (0.242) & (0.191) & (0.086) & (0.209) \\
Lab $\times$ period FE & 0.151 & $-$0.150 & 0.586 & 0.574 & 0.012 \\
 & (0.257) & (0.275) & (0.129) & (0.081) & (0.153) \\
OP-style control & $-$0.038 & $-$0.064 & 0.424 & 0.480 & $-$0.056 \\
 & (0.253) & (0.186) & (0.059) & (0.102) & (0.117) \\
\addlinespace
\multicolumn{6}{@{}l}{\textit{Panel B. Winogrande log odds (57 models)}} \\
OLS & $-$0.037 & $-$0.234 & 0.351 & 0.511 & $-$0.160 \\
 & (0.192) & (0.161) & (0.044) & (0.092) & (0.102) \\
Year FE & $-$0.001 & $-$0.265 & 0.339 & 0.499 & $-$0.161 \\
 & (0.198) & (0.173) & (0.045) & (0.084) & (0.096) \\
Developer FE & 0.105 & $-$0.297 & 0.399 & 0.498 & $-$0.099 \\
 & (0.189) & (0.193) & (0.050) & (0.078) & (0.093) \\
Family FE & 0.225 & $-$0.549 & 0.405 & 0.466 & $-$0.061 \\
 & (0.171) & (0.164) & (0.120) & (0.072) & (0.140) \\
Lab $\times$ period FE & 0.153 & $-$0.278 & 0.425 & 0.471 & $-$0.046 \\
 & (0.199) & (0.240) & (0.072) & (0.066) & (0.097) \\
OP-style control & 0.006 & $-$0.234 & 0.336 & 0.494 & $-$0.158 \\
 & (0.202) & (0.161) & (0.055) & (0.081) & (0.097) \\
\addlinespace
\multicolumn{6}{@{}l}{\textit{Panel C. HellaSwag log odds, restricted sample (49 models)}} \\
OLS & 0.138 & $-$0.013 & 0.382 & 0.337 & 0.045 \\
 & (0.079) & (0.062) & (0.021) & (0.009) & (0.023) \\
Year FE & 0.141 & $-$0.066 & 0.373 & 0.357 & 0.016 \\
 & (0.085) & (0.089) & (0.027) & (0.008) & (0.028) \\
Developer FE & 0.141 & $-$0.076 & 0.432 & 0.394 & 0.038 \\
 & (0.085) & (0.145) & (0.026) & (0.010) & (0.028) \\
Family FE & 0.228 & $-$0.344 & 0.462 & 0.426 & 0.036 \\
 & (0.093) & (0.223) & (0.067) & (0.014) & (0.068) \\
Lab $\times$ period FE & 0.156 & $-$0.090 & 0.466 & 0.418 & 0.047 \\
 & (0.096) & (0.137) & (0.045) & (0.010) & (0.046) \\
OP-style control & 0.134 & $-$0.027 & 0.391 & 0.365 & 0.026 \\
 & (0.092) & (0.065) & (0.022) & (0.007) & (0.023) \\
\bottomrule
\end{tabular}
\par\smallskip
\begin{tablenotes}Same design-matched comparison as Table~\ref{tab:lalonde}. Panels A and B: 57 models from 19 developers; Panel C excludes Qwen1.5 (token counts imputed by size), distilled, synthetic-data, and code models (49 models, 16 developers). Each estimator is applied to the observed benchmark log odds of the models inside the convex hull of the experimental runs (``Observational'') and, on the same models, to counterfactual outputs from the experimental technology, a quadratic surface fitted to the OLMo ladder and the \citet{gadre2024language} testbed (``Experimental''). Bias $=$ Observational $-$ Experimental. Standard errors in parentheses: clustered by developer for the observational estimates, from 300 bootstrap refits of the surface for the benchmark; the bias standard error treats the two as independent. With 16 to 19 developer clusters these standard errors can overstate precision; in Panel C the wild-cluster bootstrap $p$-value of the OLS $\theta_C$ bias is 0.10 (normal $p$ 0.03). $\theta_N$ and $\theta_D$ are the elasticities of the success odds with respect to parameters and tokens; the OP-style row controls for a quadratic in the estimated probability of Epoch notability.\end{tablenotes}
\begin{tablenotes}[Source]Open LLM Leaderboard evaluations via \citet{ruan2024observational} and \citet{maiapolo2024sloth}; OLMo ladder \citep{bhagia2024establishing}; \citet{gadre2024language}; authors' estimates.\end{tablenotes}
\end{table}

\begin{table}[tp]
\centering
\caption{Productivity Dispersion across Model Families in Compute-Equivalent Units}
\label{tab:app-tfp}
\footnotesize
\setlength{\tabcolsep}{2.5pt}
\begin{tabular}{@{}lcccccc@{}}
\toprule
 & & \multicolumn{2}{c}{Family effects, 90/10} & \multicolumn{2}{c}{Within developer, 90/10} & Loss \\
\cmidrule(lr){3-4}\cmidrule(lr){5-6}
Sample & Models & Exp.\ $\theta$ & Exp.\ $\theta$, $\ln C$ & Own $\theta$ & Exp.\ $\theta$ & units \\
\midrule
\multicolumn{7}{@{}l}{\textit{Panel A. HellaSwag log odds}} \\
All models & 128/38 & 23.7 & 19.9 & 11.1 & 22.4 & 1.76 \\
 & & [9.1, 197] & [7.6, 161] & [4.4, 18.8] & [6.3, 34.6] & \\
No distilled, synthetic & 121/36 & 18.3 & 17.6 & 11.3 & 23.9 & 1.68 \\
 & & [6.9, 209] & [7.4, 138] & [4.1, 18.3] & [6.6, 37.0] & \\
Also no code & 107/32 & 7.0 & 9.1 & 8.4 & 16.4 & 1.41 \\
 & & [4.7, 168] & [5.2, 11.6] & [4.1, 13.2] & [5.3, 28.6] & \\
\addlinespace
\multicolumn{7}{@{}l}{\textit{Panel B. Composite of HellaSwag, ARC-Challenge, and Winogrande}} \\
All models & 128/38 & 21.5 & 12.5 & 11.8 & 17.9 & 1.73 \\
 & & [8.3, 98.7] & [5.8, 68.6] & [5.1, 16.1] & [6.9, 21.3] & \\
No distilled, synthetic & 121/36 & 12.0 & 11.2 & 11.3 & 17.8 & 1.56 \\
 & & [4.9, 119] & [5.3, 50.7] & [4.6, 16.0] & [6.2, 22.9] & \\
Also no code & 107/32 & 6.3 & 6.9 & 9.6 & 14.4 & 1.39 \\
 & & [3.4, 81.8] & [3.5, 8.8] & [3.5, 13.8] & [3.6, 20.7] & \\
\addlinespace
\multicolumn{7}{@{}l}{\textit{Panel C. Frontier models: Epoch Capabilities Index on Epoch training compute}} \\
All & 104/23 & & & 8.1 & & \\
 & & & & [3.4, 19.2] & & \\
``Confident'' compute & 66/17 & & & 3.5 & & \\
 & & & & [2.6, 5.6] & & \\
\bottomrule
\end{tabular}
\par\smallskip
\begin{tablenotes}Models: models/families. Compute-equivalent 90/10 ratios $\exp[(\omega_{90}-\omega_{10})/\theta_C]$: the factor by which the
10th-percentile family (or model) would have to scale training compute to match the 90th percentile. Family effects
$\omega_f$ come from a regression of the output on $\ln N$, $\ln D$, and family fixed effects; the ``$\ln C$'' column
instead nets inputs with the experimental compute elasticity ($\omega_f$ = family mean of $y-\theta_C\ln C$), so that
netting and conversion use one technology. Within-developer residuals come from a regression of the output on $\ln C$
with developer and release-year effects, the design of \citet{mertens2026secret}; developers with a single model are
excluded. Own $\theta$: the regression's compute elasticity (0.43 for HellaSwag, 0.45 for the composite, 4.76 index
points per log point for the ECI); experimental $\theta$: the design-matched benchmark of Table~\ref{tab:lalonde}
(0.329 for HellaSwag, 0.388 for the composite). Brackets: 90 percent intervals from 300 developer-cluster bootstrap
draws, holding the experimental $\theta$ fixed. Loss units raise the family ratio (experimental $\theta$) to the
power $\gamma=0.178$, an output-cardinalization choice; with returns to scale near one, the 90/10 ratio of 1.92 for
US manufacturing \citep{syverson2011determines} is also an input-equivalent ratio and compares with the
compute-equivalent columns. Distilled: Gemma 2 2B and 9B; synthetic-data: Phi-1.5, Phi-2, SmolLM; code: CodeLlama,
StarCoder, StarCoder2, DeepSeek-Coder.\end{tablenotes}
\begin{tablenotes}[Source]Open LLM Leaderboard evaluations via \citet{ruan2024observational} and
\citet{maiapolo2024sloth}; Epoch AI \citep{epochai2026data,epochai2026capabilities}; authors'
estimates.\end{tablenotes}
\end{table}


\begin{figure}[tp]
\centering
\includegraphics[width=0.8\textwidth]{m4_observational_tfp_families}
\caption{Family Productivity in Compute-Equivalent Units}\label{fig:app-tfp-families}
\begin{figurenotes}
Family fixed effects from a regression of HellaSwag log odds on $\ln N$, $\ln D$, and family indicators (128 dense
base models, 38 families), converted to $\log_{10}$ compute-equivalents with the design-matched experimental compute
elasticity (0.329) and centered at the median family. The shaded region spans the 10th to 90th percentiles (a ratio of
23.7). Colors mark general, code-specialized, and distilled or synthetic-data families. Source: Open LLM Leaderboard
evaluations via \citet{ruan2024observational} and \citet{maiapolo2024sloth}; authors' estimates.
\end{figurenotes}
\end{figure}

\clearpage

\subsection{Algorithmic Progress}\label{app:add-progress}

\emph{Replication.}---Table~\ref{tab:app-ho-replication} reproduces the preferred estimates of
\citet{ho2024algorithmic} exactly: objective, parameters, and the bootstrap median of the doubling time of effective
compute (8.44 months; the point estimate is 8.68). The published point is not the minimizer of the authors' own
objective: run to convergence, the same algorithm lowers the objective by 2 percent and moves the doubling time to 6.1
months with a paper-cluster interval of $[3.0,22.7]$. Almost all of that gain comes from the penalty on parameters
whose values depend on an arbitrary normalization; the fit improves by 0.3 percent. The point estimate therefore
depends on estimator choices (6.1 to 10.2 months) that the data cannot arbitrate, and the profile-likelihood interval
is $[4.1,40.5]$ months.

\emph{Why the cross-lab data elasticity is small.}---Table~\ref{tab:app-attenuation} decomposes the gap between the
cross-lab data exponent ($\hat\beta$ of 0.025 to 0.04) and experimental values near 0.37. Dropping the irreducible
loss is the key specification choice: on the Chinchilla sweep itself, the $E=0$ form collapses one exponent and lowers
the frontier elasticity from 0.178 to 0.053, the sample-mean elasticity of total rather than reducible loss. The
like-for-like gap between cross-lab estimates (0.021 to 0.027) and experiments is therefore about 2 to 2.5 times, not 7
to 9. Because dropping $E$ rescales exponents and year coefficients proportionally to first order, the progress rate
is much less affected than the exponents when the $E=0$ model is fitted to its global optimum (pseudo-true doubling
time 11.1 months against a true 12); the authors' procedure, however, can stop at a local optimum that overstates
progress (4.6 months when $D$ is measured as dataset size). How data are measured moves the estimate most: counting
effective data under repetition raises the doubling time to 25.7 months (12.7 with the authors' stopping rule). Within
the common-$E$ form, Ho et al.'s own data prefer $E=0$, so the case for this channel rests on the sweep and the Monte
Carlo.

\emph{Architecture gains and frontier regressions.}---Table~\ref{tab:app-ceg-sahal} shows that the data cannot tell
whether the compute-equivalent gain of transformers over other architectures depends on scale: the difference in
frontier elasticities is 0.008 $[-0.012,0.034]$ and the implied gains span six orders of magnitude. Regressions of loss
on compute without a time trend overstate the compute elasticity by a factor of 1.15 on all rows and 1.44 among
record-setting releases, within the range $\gamma/(1-s_A)$ predicts for algorithmic shares between 0.05 and 0.40.

\emph{Allocative versus technical progress.}---Figure~\ref{fig:app-allocative} decomposes the Kaplan-to-Chinchilla
rebalancing. Among models trained with at least $10^{23}$ FLOP, the geometric-mean Farrell cost efficiency rose by a
factor of 1.05 between 2020--21 and 2022--24 under the Besiroglu technology, 1.91 and 2.41 under the two versions of
Hoffmann's, and 1.58 under a Farseer fit: 2 to 40 percent of the gain implied by the Ho et al.\ rate over the same
interval. Across the other sweep technologies we estimate, the realized gain ranges from 0.55 to 4.3, and
Meta's own technology implies a loss (0.74), so even its sign depends on the technology. The 10-fold gain reported by
\citet{gundlach2025origin} is instead the counterfactual cost of continuing Kaplan's allocation rule, which reaches 5.2
at $3.8\times10^{25}$ FLOP and 11.7 at $5\times10^{26}$ FLOP under the Besiroglu technology.

\begin{table}[tp]
\centering
\caption{Replication of \citet{ho2024algorithmic}: Published, Replicated, and Converged Estimates}
\label{tab:app-ho-replication}
\footnotesize
\setlength{\tabcolsep}{5pt}
\begin{tabular}{@{}lccc@{}}
\toprule
 & Published & Replication, & Same objective, \\
 & estimates & authors' code & run to convergence \\
\midrule
$\alpha_{\text{const}}$ & 0.913 & 0.913 & 0.000 \\
 & [0.000, 1.208] & [0.000, 1.268] & [0.000, 1.373] \\
$\alpha_{\text{year}}$ & 0.004 & 0.004 & $-$0.059 \\
 & [$-$0.058, 0.032] & [$-$0.070, 0.038] & [$-$0.078, 0.045] \\
$\alpha_{\text{param}}$ & 0.068 & 0.068 & 0.138 \\
 & [0.045, 0.127] & [0.041, 0.146] & [0.042, 0.188] \\
$\beta_{\text{const}}$ & 0.771 & 0.771 & 1.295 \\
 & [0.233, 1.293] & [0.265, 1.340] & [0.000, 1.388] \\
$\beta_{\text{year}}$ & 0.036 & 0.036 & 0.045 \\
 & [$-$0.002, 0.080] & [$-$0.006, 0.084] & [$-$0.022, 0.074] \\
$\beta_{\text{data}}$ & 0.040 & 0.040 & 0.025 \\
 & [0.023, 0.062] & [0.020, 0.074] & [0.017, 0.148] \\
Doubling time $T_C$ (months) & 8.44 & 8.68 & 6.08 \\
 & [4.52, 14.27] & [4.07, 16.93] & [3.05, 22.67] \\
\addlinespace
$T_C$, bootstrap median & 8.44 & 8.58 & 6.77 \\
corr$(\alpha_{\text{year}},\beta_{\text{year}})$ across draws & & $-$0.73 & $-$0.82 \\
Objective & 0.05181 & 0.05181 & 0.05072 \\
\quad Mean squared error & & 0.04642 & 0.04629 \\
\quad Sum of absolute parameters (weight 0.0025) & & 2.159 & 1.773 \\
\bottomrule
\end{tabular}
\par\smallskip
\begin{tablenotes}Preferred model (model 7) of \citet{ho2024algorithmic}: log perplexity
$=\exp(\alpha^b_{\text{const}}-\alpha_{\text{year}}(t-t_0)-\alpha_{\text{param}}\ln(N/N_0))+\exp(\beta^b_{\text{const}}
-\beta_{\text{year}}(t-t_0)-\beta_{\text{data}}\ln(D/D_0))$ on WikiText-103, Penn Treebank, and WikiText-2
(benchmark-specific constants for the latter two not shown), with no irreducible loss; 231 model--benchmark rows from
144 papers. Estimator: mean squared error plus $0.0025\sum|\theta|$, minimized by SLSQP from $\theta=0$. Our
replication of the authors' pipeline and code matches their objective (0.0518129) and parameters to
$4\times10^{-7}$; with their iid bootstrap protocol we obtain a median $T_C$ of 8.44 and an interval of
$[4.39,14.05]$. The last column runs the same objective, algorithm, and start with a tight tolerance instead of the
default; the objective falls by 2 percent, almost entirely through the penalty term, and 300 random starts all end
below the published objective. $T_C$ is the doubling time of effective compute; the published 8.4 months is a
bootstrap median. Brackets: published intervals and, for our estimates, paper-cluster bootstrap intervals (400 draws)
that map percentiles of the growth rate (4 of 400 converged draws have nonpositive growth). The profile-likelihood 95
percent interval for $T_C$ is $[4.1,40.5]$ months.\end{tablenotes}
\begin{tablenotes}[Source]Data and code of \citet{ho2024algorithmic}; authors' estimates.\end{tablenotes}
\end{table}

\begin{table}[tp]
\centering
\caption{Why Is the Cross-Lab Data Elasticity Small? Channels of Attenuation}
\label{tab:app-attenuation}
\footnotesize
\setlength{\tabcolsep}{3.5pt}
\begin{tabular}{@{}lccccccc@{}}
\toprule
 & & & & & & \multicolumn{2}{c}{Default stop} \\
\cmidrule(l){7-8}
 & $n$ & $\hat\alpha$ & $\hat\beta$ & $\hat\gamma$ & $T_C$ [95\% CI] & $\hat\beta$ & $T_C$ \\
\midrule
\multicolumn{8}{@{}l}{\textit{Panel A. Cross-lab model on alternative data}} \\
Baseline (model 7) & 231 & 0.138 & 0.025 & 0.021 & 6.1 [3.0, 26.5] & 0.040 & 8.7 \\
Benchmark-specific perplexity & 231 & 0.139 & 0.023 & 0.020 & 5.7 [2.8, 21.2] & 0.030 & 9.0 \\
(a) Tokens seen, epochs $=1$ & 231 & 0.065 & 0.035 & 0.023 & 8.0 [1.1, $\infty$] & 0.018 & 2.6 \\
(a) Tokens seen, epochs imputed & 231 & 0.121 & 0.028 & 0.023 & 4.7 [2.4, $\infty$] & 0.027 & 4.7 \\
(a) Effective data & 231 & 0.057 & 0.067 & 0.031 & 25.7 [3.9, $\infty$] & 0.052 & 12.7 \\
(a) Word-level vocabulary & 134 & 0.040 & 0.000 & 0.000 & n.d. & 0.008 & 0.9 \\
(a) Vocabulary control & 185 & 0.135 & 0.021 & 0.018 & 3.2 [1.4, 18.7] & 0.036 & 4.6 \\
(b) Published 2018 or later & 180 & 0.052 & 0.077 & 0.031 & 7.8 [2.0, $\infty$] & 0.040 & 4.3 \\
(b) Transformers only & 104 & 0.061 & 0.039 & 0.024 & 5.2 [1.6, $\infty$] & 0.032 & 11.3 \\
(b) No year terms & 231 & 0.063 & 0.084 & 0.036 & -- & 0.037 & -- \\
(c) Benchmark-specific $E$ & 231 & 0.072 & 0.042 & 0.026 & 10.5 [2.7, 44.3] & 0.042 & 10.5 \\
(d) All models per paper & 261 & 0.055 & 0.064 & 0.030 & 15.3 [2.7, 79.3] & 0.038 & 7.6 \\
(d) Top model per paper & 144 & 0.143 & 0.025 & 0.022 & 6.7 [2.5, 56.9] & 0.042 & 9.6 \\
(d) Flagged outliers kept & 238 & 0.150 & 0.023 & 0.020 & 5.2 [2.5, 17.9] & 0.030 & 7.5 \\
(d) No zero-shot evaluations & 207 & 0.049 & 0.091 & 0.032 & 16.6 [3.1, 23.5] & 0.059 & 11.4 \\
\addlinespace
\multicolumn{8}{@{}l}{\textit{Panel B. Chinchilla sweep}} \\
$E$ estimated, Huber & 240 & 0.347 & 0.367 & \multicolumn{2}{l}{0.178 [0.169, 0.188]} &  &  \\
$E=0$, Huber & 240 & 0.375 & 0.061 & \multicolumn{2}{l}{0.053 [0.052, 0.054]} &  &  \\
$E=0$, least squares & 240 & 0.061 & 0.377 & \multicolumn{2}{l}{0.053} &  &  \\
Total-loss elasticities & 240 & 0.053 & 0.056 & \multicolumn{2}{l}{0.054} &  &  \\
\addlinespace
\multicolumn{8}{@{}l}{\textit{Panel C. Monte Carlo on the cross-lab design}} \\
 & & \multicolumn{3}{c}{Global optimum} & MC & \multicolumn{2}{c}{Authors' code} \\
\cmidrule(lr){3-5}\cmidrule(lr){6-6}\cmidrule(l){7-8}
 & & $\hat\alpha$ & $\hat\beta$ & $\hat\gamma$ & $T_C$ [5, 95] & $\hat\beta$ & $T_C$ \\
Correct specification &  & 0.348 & 0.366 & 0.178 & 9.9 [7.4, 12.8] & 0.366 & 12.0 \\
(c) $E$ dropped &  & 0.046 & 0.388 & 0.041 & 11.2 [8.8, 15.5] & 0.346 & 11.1 \\
(c) and neutrality imposed &  & 0.046 & 0.489 & 0.042 & 4.4 [3.7, 5.4] & 0.446 & 4.2 \\
(c)+(a) $D$ = dataset size &  & 0.032 & 0.389 & 0.029 & 10.4 [7.9, 16.5] & 0.048 & 4.6 \\
(c)+(a)+(d) best 3 &  & 0.036 & 0.415 & 0.033 & 10.9 [7.6, 18.1] & 0.280 & 11.9 \\
(c)+(a) random 3 &  & 0.033 & 0.398 & 0.030 & 10.7 [7.5, 17.9] & 0.303 & 10.5 \\
(c)+(a), $\rho(\psi_D,\ln D)=-0.5$ &  & 0.031 & 0.370 & 0.028 & 9.8 [5.3, 16.5] & 0.045 & 4.6 \\
(c)+(a), $\rho(\psi_D,\ln D)=+0.5$ &  & 0.025 & 0.384 & 0.023 & 10.5 [5.8, 21.4] & 0.288 & 10.1 \\
\bottomrule
\end{tabular}
\par\smallskip
\begin{tablenotes}Panel A: the model and penalized objective of \citet{ho2024algorithmic} (mean squared error plus $0.0025\sum|\theta|$, $E=0$), minimized to convergence on the indicated data; the last two columns (default stop) stop the optimizer at the default tolerance used in the authors' code. $T_C$ is the doubling time of effective compute, with a 95 percent paper-cluster bootstrap interval (400 draws, each started at zero; ``$\infty$'' when nonpositive growth is inside the interval). Rows: (a) measurement of data, with $D$ as tokens seen (missing epoch counts set to 1 or imputed from dataset size), as effective data under repetition \citep[$R^*=15.4$;][]{muennighoff2023scaling}, or with word-level vocabularies only or a vocabulary control; 33 percent of epoch counts are imputed; (b) time and collinearity; (c) benchmark-specific irreducible loss estimated by NLS; (d) selection, using all models of each paper, only the best, keeping flagged outliers, or dropping zero-shot evaluations. Panel B: the \citet{besiroglu2024chinchilla} estimator on the Chinchilla sweep with $E$ estimated or set to zero, with 95 percent bootstrap intervals for $\gamma=\alpha\beta/(\alpha+\beta)$; the last row gives the sample means of the total-loss elasticities $\alpha u/L$, $\beta v/L$, and $\gamma(L-E)/L$ implied by the full model. Panel C: pseudo-true values (at the global optimum and with the authors' code) from noise-free data when the truth is the Besiroglu technology with Hicks-neutral progress ($T_C=12$ months) evaluated at the actual inputs of the sample; MC median and 5th--95th percentiles of $T_C$ from 204 replications with noise of standard deviation 0.215; rows switch on (c) dropping $E$, (a) measuring $D$ as dataset size, (d) keeping the best three models per paper, and a data-augmenting productivity shock $\psi_D$ correlated with $\ln D$. n.d.: not defined because $\hat\beta=0$.\end{tablenotes}
\begin{tablenotes}[Source]Data of \citet{ho2024algorithmic}; Epoch AI digitization of \citet{hoffmann2022training}; authors' estimates.\end{tablenotes}
\end{table}

\input{tables/appD_ceg_sahal}

\begin{figure}[tp]
\centering
\includegraphics[width=\textwidth]{m5_progress_fig6_allocative}
\caption{Allocative Efficiency and the Kaplan-to-Chinchilla Rebalancing}\label{fig:app-allocative}
\begin{figurenotes}
Panel (a): Farrell cost efficiency $C_{\min}(L)/6ND$ of 247 Epoch language models under the
\citet{besiroglu2024chinchilla} technology, with era geometric means for models trained with at least $10^{23}$ FLOP;
efficiency is measured against the training-only frontier, so models over-trained for inference score below 1 by
design. Panel (b): compute-equivalent cost of Kaplan's allocation rule ($N\propto C^{0.73}$, anchored at GPT-3) under
three technologies, with the 10-fold gain of \citet{gundlach2025origin}. Panel (c): allocative gain between 2020--21
and 2022--24 as a share of the gain implied by the effective-compute growth rate of \citet{ho2024algorithmic}, by
technology and sample. $D$ is inferred as $C/(6N)$. Source: Epoch AI \citep{epochai2026data}; authors' estimates.
\end{figurenotes}
\end{figure}

\clearpage

\subsection{Bias Formulas for Input Measurement and Flexible Inputs}\label{app:add-biasformulas}

This subsection derives the two bias formulas used in Section~\ref{sec:tech:measurement}. Both were verified
symbolically and against brute-force IsoFLOP argmins.

\emph{Input measurement.}---Let the true technology be $L=E+AN^{-\alpha}+BD^{-\beta}$ and let the researcher count
$N_m$ parameters, where the true count is $N=N_m+X(N_m)$ and the omitted component has share $s\equiv X/N$ and a
constant scale elasticity $\theta\equiv d\ln X/d\ln N_m<1$. The researcher measures compute as $C_m=6N_mD$ and takes
the IsoFLOP argmin in measured units: $n_m\equiv\ln N_m$ minimizes $L$ subject to $\ln D=c_m-\ln6-n_m$, with
$c_m\equiv\ln C_m$. Because $d\ln N/dn_m=(N_m+\theta X)/N=1-s(1-\theta)\equiv\varphi$, the first-order condition is
$\alpha u\varphi=\beta v$, with $u=AN^{-\alpha}$ and $v=BD^{-\beta}$. Differentiating its log along the researcher's
path, with $d\ln u/dn_m=-\alpha\varphi$, $d\ln v=\beta\,(dn_m-dc_m)$ and $\zeta\equiv d\ln\varphi/dn_m$, gives
\[
(\zeta-\alpha\varphi)\,\frac{dn_m}{dc_m}=\beta\,\frac{dn_m}{dc_m}-\beta
\qquad\Longrightarrow\qquad
a_m\equiv\frac{dn_m}{dc_m}=\frac{\beta}{\beta+\alpha\varphi-\zeta},
\]
which is equation~\eqref{eq:measbias}. With $\theta$ constant, $d\ln s/dn_m=\theta-\varphi$, so
$\zeta=(1-\theta)s(\varphi-\theta)/\varphi\ge0$, because $\varphi\ge\theta$. Since $\varphi\le1$ and $\zeta\ge0$, the
denominator is at most $\alpha+\beta$, and $a_m\ge a=\beta/(\alpha+\beta)$, with equality when $s=0$ or $\theta=1$. A
proportional error ($\theta=1$) gives $\varphi=1$ and $\zeta=0$ and moves only the level of the path. As $s\to1$,
$a_m\to\beta/(\beta+\alpha\theta)$, the small-scale limit of \citet{pearce2024reconciling}. The closed form matches
brute-force argmins to within $2\times10^{-7}$ in $\ln N$.

\emph{Flexible inputs.}---Let observed loss be $L_{\rm obs}=L^*e^{\iota(N,D)}$, where $L^*$ is the technology with
the hyperparameters concentrated out and $\iota\ge0$ is the inefficiency of the researcher's hyperparameter policy.
Along the IsoFLOP curve at compute $C$, write $f(n)$ for $L^*$ at $d=c-\ln6-n$, with minimizer $n^*(c)$; then
$f'=-\alpha u+\beta v$ and $f''=\alpha^2u+\beta^2v$, and at the minimizer $\alpha u=\beta v=\gamma R^*$, with $R^*$ the
reducible loss on the frontier, so $f''=(\alpha+\beta)\gamma R^*$. The observed argmin solves $f'(n)+\Delta(n)=0$,
where $\Delta\equiv\partial(L_{\rm obs}-L^*)/\partial n$ at fixed compute. Expanding around $n^*(c)$, the log argmin
shifts by $-\Delta(C)/f''(C)$ to first order in $\Delta$, and differentiating in $\ln C$ gives
$a_{\rm obs}-a=-\partial[\Delta/f'']/\partial\ln C$. At the minimizer $\Delta=L^*e^{\iota}(\iota_n-\iota_d)$, with
$\iota_n\equiv\partial\iota/\partial\ln N$ and $\iota_d\equiv\partial\iota/\partial\ln D$, so inefficiency that depends
on the inputs only through compute ($\iota_n=\iota_d$) leaves $a$ unbiased. With small, constant log-gradients,
$L^*=E+R^*$ and $R^*\propto C^{-\gamma}$ give $a_{\rm obs}-a=-(\iota_n-\iota_d)E/[(\alpha+\beta)R^*(C)]$. Under the
\citet{besiroglu2024chinchilla} technology the formula matches simulated IsoFLOP argmins to within 0.003 in $a$.

\subsection{Hyperparameters as Flexible Inputs}\label{app:add-flexible}

Table~\ref{tab:measurement} decomposes the Kaplan--Chinchilla discrepancy into input measurement and flexible inputs.
The tables below use the Step Law grid of \citet{li2025predictablea}, which trains every $(N,D)$ cell under many
learning rates and batch sizes, to study hyperparameters directly. Table~\ref{tab:app-steplaw} shows that this
17-cell design does not identify the allocation exponent: the likelihood is flat in $E$ while $a$ ranges from 0.76 to
0.45. Holding $E$ fixed, published hyperparameter rules shift $a$ by $-0.041$ to $+0.076$, all within 0.7
residual-bootstrap standard errors. Choosing one random configuration per cell costs mainly precision for $a$ (median
bias $-0.3$ standard errors) but biases $\sigma^*$ ($-0.051$, or $-3.7$ standard errors), so these data give no reason
to regard $\sigma^*$ as the more robust object. The stochastic frontier in Table~\ref{tab:app-sfa} shows that
inefficiency falls with training tokens ($\gamma_D$ between $-0.23$ and $-0.26$) but not significantly with
parameters: longer horizons flatten the loss surface in the hyperparameters. Because only the transverse gradient of
inefficiency biases the IsoFLOP exponent, this pattern implies a small downward bias ($-0.012$) for a researcher who
uses random configurations. Table~\ref{tab:app-demand} estimates the demand for the flexible inputs themselves. The
unconditional learning-rate demand replicates Step Law's; holding the batch fixed, the positive token elasticity
shrinks to between 0.02 and 0.05, but conditioning closes only 26 to 43 percent of the gap to the negative elasticity
reported by \citet{bjorck2024scaling}.

\begin{table}[tp]
\centering
\caption{Hyperparameters as Flexible Inputs: Inefficiency and the Estimated Technology}
\label{tab:app-steplaw}
\footnotesize
\setlength{\tabcolsep}{2.5pt}
\begin{tabular}{@{}lccccccc@{}}
\toprule
 & \multicolumn{3}{c}{Inefficiency $\iota$ (percent)} & \multicolumn{2}{c}{$E$ held at 1.4} & \multicolumn{2}{c}{$E$ estimated} \\
\cmidrule(lr){2-4}\cmidrule(lr){5-6}\cmidrule(l){7-8}
Policy & Mean & $\partial\iota/\partial\ln N$ & $\partial\iota/\partial\ln D$ & $a$ & $\sigma^*$ & $a$ & $\sigma^*$ \\
\midrule
Frontier (best run) & 0.00 & 0.00 & 0.00 & 0.573 & 0.790 & 0.563 & 0.786 \\
 & & & & (0.108) & (0.014) & (0.148) & (0.040) \\
Step Law rule & 0.10 & $-$0.05 (0.04) & 0.01 (0.02) & 0.532 & 0.794 & 0.517 & 0.789 \\
Best of 16 random & 0.21 & $-$0.04 [0.08] & $-$0.06 [0.06] & 0.558 & 0.791 & 0.551 & 0.788 \\
Best of 4 random & 0.70 & $-$0.16 [0.25] & $-$0.20 [0.20] & 0.555 & 0.786 & 0.589 & 0.793 \\
One random & 2.58 & $-$0.44 [1.19] & $-$0.70 [0.71] & 0.539 & 0.739 & 0.565 & 0.760 \\
One fixed setting & 0.29 & $-$0.08 (0.08) & $-$0.03 (0.05) & 0.619 & 0.780 & 0.515 & 0.748 \\
Porian et al., base & 0.41 & 0.32 (0.16) & $-$0.23 (0.11) & 0.621 & 0.785 & 0.564 & 0.769 \\
Porian et al., $N$-rule & 0.77 & 0.72 (0.22) & $-$0.57 (0.14) & 0.649 & 0.785 & 0.453 & 0.722 \\
DeepSeek $C$-rule & 0.75 & $-$0.16 (0.06) & $-$0.02 (0.03) & 0.543 & 0.793 & 0.488 & 0.774 \\
Bj\"orck et al.\ rule & 0.91 & $-$0.57 (0.07) & 0.04 (0.04) & 0.576 & 0.782 & 0.548 & 0.773 \\
\bottomrule
\end{tabular}
\par\smallskip
\begin{tablenotes}The Step Law grid of \citet{li2025predictablea}: 1,911 runs in 17 $(N,D)$ cells with $N$ from 215
million to 1.07 billion non-embedding parameters and $D$ from 4 to 100 billion tokens, up to 12 learning rates and 10
batch sizes per cell; the output is the smoothed final training loss. Each policy selects one run per cell (the grid
point nearest its target); random rows are medians over 200 Monte Carlo draws that keep the best of $k$ non-diverged
configurations per cell (Monte Carlo standard deviations in brackets). Inefficiency gradients: OLS across the 17 cells,
HC1 standard errors in parentheses. Technology: the Chinchilla form fitted by NLS on log loss, with residual-bootstrap
standard errors for the frontier (200 draws). The profile likelihood in $E$ is flat (sum of squared residuals within
17 percent for $E$ between 0.6 and 1.7, over which $a$ ranges from 0.76 to 0.45), so only differences across rows at
fixed $E$ are informative. $N$ follows Step Law's non-embedding convention, which biases the level of $a$ upward; counting
the output head moves the frontier $a$ to 0.465 ($E$ estimated) and 0.550 ($E=1.4$). The Huber estimator on the frontier
gives $E=0.91$, $a=0.689$, and $\sigma^*=0.824$. $\iota=\ln L-\ln L^*$. ``One fixed setting'' uses the best single (learning rate, batch) pair for all cells. Rules: Step Law, the learning-rate and batch rules of
\citet{li2025predictablea}; DeepSeek, the compute-based rules of \citet{bi2024deepseek}; Bj\"orck et al., the learning-rate rule
in $(N,D)$ of \citet{bjorck2024scaling} at a fixed 0.5-million-token batch; Porian et al., the base setting and the
$N$-based rule of \citet{porian2024resolving}.\end{tablenotes}
\begin{tablenotes}[Source]\citet{li2025predictablea}; authors' estimates.\end{tablenotes}
\end{table}

\begin{table}[tp]
\centering
\caption{Stochastic Frontier with Scale-Dependent Inefficiency}
\label{tab:app-sfa}
\footnotesize
\setlength{\tabcolsep}{4pt}
\begin{tabular}{@{}llcccc@{}}
\toprule
Inefficiency & Noise s.d.\ $\sigma_v$ & $\gamma_0$ & $\gamma_N$ & $\gamma_D$ & Log likelihood \\
\midrule
Half-normal & Estimated & $-$3.31 & $-$0.038 & $-$0.231 & 4,462.8 \\
 & & (0.02) [0.03] & (0.037) [0.076] & (0.020) [0.032] & \\
Half-normal & Calibrated, 0.0017 & $-$3.31 & $-$0.014 & $-$0.229 & 4,423.5 \\
 & & (0.02) & (0.038) & (0.020) & \\
Exponential & Estimated & $-$3.68 & $-$0.063 & $-$0.260 & 4,585.0 \\
 & & (0.03) [0.03] & (0.053) [0.070] & (0.028) [0.023] & \\
Exponential & Calibrated, 0.0017 & $-$3.70 & $-$0.056 & $-$0.262 & 4,575.7 \\
 & & (0.02) & (0.054) & (0.028) & \\
\addlinespace
\multicolumn{2}{@{}l}{Cell-level OLS of $\ln\bar\iota_c$} & $-$3.68 & $-$0.134 & $-$0.266 & \\
 & & (0.04) & (0.108) & (0.028) & \\
\bottomrule
\end{tabular}
\par\smallskip
\begin{tablenotes}Cost-type stochastic frontier $\ln L_{ij}=\mu_c+\epsilon_{ij}+\iota_{ij}$ on the 1,730 non-diverged
runs of the Step Law grid \citep{li2025predictablea}, with 17 cell effects $\mu_c$ (the concentrated technology, left
unrestricted), noise $\epsilon\sim N(0,\sigma_v^2)$, and inefficiency $\iota\geq0$ that is half-normal or exponential
with $\ln\sigma_{u}=\gamma_0+\gamma_N(\ln N-\overline{\ln N})+\gamma_D(\ln D-\overline{\ln D})$ (heteroskedastic
inefficiency, \citealt{caudill1995frontier}). Runs with loss more than 35 percent above the cell minimum (9.5 percent of
runs) are treated as diverged. Calibrated $\sigma_v$ is the median residual standard deviation of the within-cell
quadratic surfaces. Maximum likelihood; standard errors from the numerical Hessian in parentheses and cell-clustered
sandwich standard errors (17 clusters) in brackets. The estimated frontier coincides with the best run in each cell
(mean absolute gap at most 0.11 percent): because the grid is designed, its minimum measures the frontier directly, and
the stochastic frontier adds only the variance function of inefficiency. Last row: OLS across the 17 cells of the log mean inefficiency on centered
$\ln N$ and $\ln D$ (HC1 standard errors), comparable with $(\gamma_N,\gamma_D)$ for exponential
inefficiency.\end{tablenotes}
\begin{tablenotes}[Source]\citet{li2025predictablea}; authors' estimates.\end{tablenotes}
\end{table}

\begin{table}[tp]
\centering
\caption{Flexible-Input Demand Functions: Learning Rate and Batch Size}
\label{tab:app-demand}
\footnotesize
\setlength{\tabcolsep}{4pt}
\begin{tabular}{@{}lcccc@{}}
\toprule
Specification & $\ln N$ & $\ln D$ & $\ln B$ & Observations \\
\midrule
\multicolumn{5}{@{}l}{\textit{Panel A. Learning-rate demand, $\ln\eta^*$}} \\
Unconditional, grid argmin & $-$0.823 (0.127) & 0.288 (0.071) & & 17 \\
Unconditional, smoothed argmin & $-$0.693 (0.050) & 0.210 (0.025) & & 17 \\
Conditional on batch, smoothed & $-$0.534 (0.027) & 0.046 (0.015) & 0.320 (0.034) & 165 \\
Conditional on batch, grid argmin & $-$0.612 (0.048) & 0.016 (0.029) & 0.295 (0.038) & 157 \\
Batch fixed at 0.52M tokens & $-$0.495 (0.077) & 0.040 (0.047) & & 16 \\
\addlinespace
\multicolumn{5}{@{}l}{\textit{Panel B. Batch-size demand, $\ln B^*$ (tokens)}} \\
Grid argmin & $-$0.355 (0.176) & 0.543 (0.075) & & 17 \\
Smoothed argmin & $-$0.205 (0.041) & 0.598 (0.021) & & 17 \\
Smoothed argmin, $D$ only & & 0.572 (0.037) & & 17 \\
\addlinespace
\multicolumn{5}{@{}l}{\textit{Panel C. Published rules}} \\
Step Law, $\eta$ & $-$0.713 (0.030) & 0.307 (0.019) & & \\
Step Law, $B$ & & 0.571 (0.019) & & \\
\citet{bjorck2024scaling}, $\eta$, fixed batch & $-$0.23 & $-$0.32 & & \\
\addlinespace
\multicolumn{5}{@{}l}{\textit{Panel D. Decomposition of $d\ln\eta^*/d\ln D$}} \\
 & Smoothed & Grid & & \\
\cmidrule(lr){2-3}
Conditional $+$ complementarity $\times$ batch & 0.238 & 0.176 & & \\
Unconditional & 0.210 & 0.288 & & \\
Share due to batch co-scaling & 0.91 (0.15) & 0.56 (0.24) & & \\
Share of published gap closed & 0.26 (0.02) & 0.43 (0.13) & & \\
\bottomrule
\end{tabular}
\par\smallskip
\begin{tablenotes}Optimal learning rate $\eta^*$ and batch size $B^*$ in each of the 17 Step Law $(N,D)$ cells, regressed
on $\ln N$ (non-embedding) and $\ln D$. ``Grid argmin'' uses the best run in the cell; ``smoothed'' uses the minimum of a
quadratic in $(\ln\eta,\ln B)$ fitted to runs within 2 percent of the cell minimum, which lies above the best observed
run in 16 of 17 cells. Standard errors in parentheses: a two-level bootstrap (cells, then each cell's optimum; 1,000
draws) for the unconditional demands and cell-clustered standard errors for the conditional ones (cell $\times$ batch
slices with an interior optimum); Panel C reports the standard deviations of Step Law's own bootstrap fits. Conditional
demands hold the batch fixed, as in \citet{bjorck2024scaling}; unconditional demands let it adjust. Panel D: the published gap is between Step Law's $+0.307$ and
Bj\"orck et al.'s $-0.32$, their joint fit for models of at least 760 million parameters; their per-size exponents are
steeper for smaller models (down to $-0.70$ at 50 million), and against $-0.70$ the shares closed would be about 0.16
(smoothed) and 0.27 (grid); the unconditional elasticity equals the conditional one plus the learning-rate--batch complementarity times the batch
elasticity (the Le Chatelier decomposition of \citealt{milgrom1996lechatelier}); cell-cluster bootstrap standard errors
in parentheses (1,000 draws).\end{tablenotes}
\begin{tablenotes}[Source]\citet{li2025predictablea}; authors' estimates.\end{tablenotes}
\end{table}


\clearpage

\subsection{Verification of the Formal Results}\label{app:add-claims}

Table~\ref{tab:app-claims} lists the formal statements of the paper with the checks applied to each. Every algebraic
identity was reduced to zero symbolically, and every comparative-static, bias, and information statement was also
checked by a method that does not use the closed form: brute-force optimization, simulation, or direct computation of
Fisher information. Three statements differ from formulations that are common in earlier drafts or readings of the
literature, and the paper uses the corrected versions: the sign of biased technical change is identified from the drift
of compute-optimal tokens per parameter, compute-equivalent gains are constant when the irreducible loss and the
frontier elasticity coincide (not the individual exponents), and forward and reverse regressions bracket the returns to
compute only when compute does not rise with productivity.

\begin{table}[tp]
\centering
\caption{Verification Register for the Formal Results}
\label{tab:app-claims}
\footnotesize
\setlength{\tabcolsep}{3pt}
\renewcommand{\arraystretch}{1.0}
\begin{tabular}{@{}>{\raggedright\arraybackslash}p{0.52\textwidth}>{\raggedright\arraybackslash}p{0.15\textwidth}>{\raggedright\arraybackslash}p{0.07\textwidth}>{\raggedright\arraybackslash}p{0.2\textwidth}@{}}
\toprule
Statement & Result & Checks & Status \\
\midrule
\multicolumn{4}{@{}l}{\textit{Technology and elasticities}} \\
$y=-\ln(L-E)=\omega+F(n+\psi_N,d+\psi_D)-\epsilon$ exactly, given $E$; $\varepsilon_N=\alpha u/(u+v)$ & Section~\ref{sec:framework} & S & Exact, not approximate \\
$\sigma=(\alpha u+\beta v)/[\alpha u(1+\beta)+\beta v(1+\alpha)]$; $\sigma^*=2/(2+\alpha+\beta)$ & Lemma~\ref{lemma:sigma} & S, N & Verified \\
$\sigma(w)=(1+w)/[1+\alpha+w(1+\beta)]$; $\sigma$ is invariant to monotone transformations of output & Lemma~\ref{lemma:sigma} & S & New \\
An interior compute optimum under $C=6ND$ is a strict minimum if and only if $0<\sigma<1$ & Lemma~\ref{lemma:sigma} & S, N & New \\
\addlinespace
\multicolumn{4}{@{}l}{\textit{Duality, allocation, and geometry}} \\
$N^*=G(C/6)^a$, $L^*(C)=E+K(C/6)^{-\gamma}$, $C^*(R)=6(K/R)^{1/\gamma}$ & Proposition~\ref{prop:duality} & S, N & Verified \\
Hicks-neutral $\omega$ does not move the mix at given compute; better data lower $M^*$ & Proposition~\ref{prop:duality} & S, N & Holds at given compute, not at a given loss \\
The Hessian of $\ln R$ in $(n,d)$ has rank one with null vector $(\beta,\alpha)$; the transverse coordinate is $-\ln w$ & Proposition~\ref{prop:fdep} & S & Verified; decomposition new \\
\addlinespace
\multicolumn{4}{@{}l}{\textit{Identification}} \\
On-path Jacobian has rank 3; $\sigma^*$ first-order estimable only if $\alpha=\beta$ & Proposition~\ref{prop:fdep} & N & Verified \\
With $\kappa$ free, every $\sigma^*\in(0,1)$ fits on-path data exactly & Proposition~\ref{prop:fdep} & S, N & Verified \\
$y=y^*(c)-\gamma\Phi-\epsilon$; information about $\ln M^*$ scales as $v^2$, about $\alpha+\beta$ as $v^4$ & Proposition~\ref{prop:info} & S, N & New \\
Time series of compute-optimal models identify $g_N$, $g_D$, and the sign of the Hicks bias, not $\sigma^*$ or the bias's magnitude & Proposition~\ref{prop:dmr} & S, N & Corrected: the sign is identified \\
Frontier compute-equivalent gains are constant in $C$ if and only if $E$ and $\gamma$ are equal & Proposition~\ref{prop:dmr} & S, N & Corrected: equal exponents are not required \\

\bottomrule
\end{tabular}
\par\smallskip
\begin{tablenotes}[Note]The register continues in the next table, which also gives the notes.\end{tablenotes}
\end{table}

\begin{table}[tp]
\addtocounter{table}{-1}\renewcommand{\theHtable}{\thetable-continued}
\centering
\caption{Verification Register for the Formal Results (continued)}
\footnotesize
\setlength{\tabcolsep}{3pt}
\renewcommand{\arraystretch}{1.0}
\begin{tabular}{@{}>{\raggedright\arraybackslash}p{0.52\textwidth}>{\raggedright\arraybackslash}p{0.15\textwidth}>{\raggedright\arraybackslash}p{0.07\textwidth}>{\raggedright\arraybackslash}p{0.2\textwidth}@{}}
\toprule
Statement & Result & Checks & Status \\
\midrule
\multicolumn{4}{@{}l}{\textit{Observational data}} \\
OLS slope of $y$ on $c$ is $\gamma+\pi_1V_\omega/(\pi_1^2V_\omega+V_\eta)$ for any sign of $\pi_1$ & Proposition~\ref{prop:transmission} & S, MC & Verified \\
Forward and reverse regressions bracket $\gamma$ if and only if $-(1+V_\epsilon/V_\omega)/\gamma\leq\pi_1\leq0$ & Proposition~\ref{prop:transmission} & S & Corrected: fails under funding \\
Hicks-neutral transmission biases $\gamma$ but not the allocation exponent $a$ & Proposition~\ref{prop:transmission} & MC & Verified \\
Outcome-based release: slope $\leq\gamma$; in $[0,\gamma]$ under log-concavity; net sign under funding flips at a closed-form threshold & Proposition~\ref{prop:transmission} & MC & $\mathrm{E}[\omega\mid c,\text{released}]$ need not fall in $c$ \\
Given $c$, $\ln(D/N)$ inverts $\chi=\beta\psi_D-\alpha\psi_N$ but is flat in $\omega$ & Proposition~\ref{prop:transmission} & S, N & New \\
\addlinespace
\multicolumn{4}{@{}l}{\textit{Revealed inference demand}} \\
$w=1+T/(3D)$, invariant to $\omega$, $\epsilon$, $E$, and $\kappa$; $\hat w=w\,e^{-\chi}$ & Proposition~\ref{prop:wedge} & S, N & Verified \\
A binding data constraint lowers $w$ below $1+T/(3D)$; a binding memory constraint raises it & Proposition~\ref{prop:wedge} & S, N & Verified \\
$w=(M/M^*(C))^{1/\sigma^*-1}$ and $C/C_{\min}=((\alpha+\beta w)/(\alpha+\beta))^{1/\gamma}w^{-1/\alpha}$ for any $\alpha,\beta$ & Proposition~\ref{prop:wedge} & S, N & New; generalizes $\alpha=\beta$ \\
$\mathrm{Var}(\ln\hat w)$ is quadratic in $\ln M$; monotone lower bound on $T$ under concavity and supermodularity & Proposition~\ref{prop:pi} & N & New \\
\addlinespace
\multicolumn{4}{@{}l}{\textit{Growth accounting}} \\
Frontier-release exponent $\gamma[1+(g_A/g)R^2_{c,t}]$, $=\gamma/(1-s_A)$ with deterministic trends & Appendix~\ref{app:proofs} & S, MC & Verified \\
Hall-type bias of equal-weight residuals, $(\varepsilon_N-\varepsilon_D)(\Delta n-\Delta d)/2$ & Appendix~\ref{app:proofs} & S, N & Verified to first order \\
\bottomrule
\end{tabular}
\par\smallskip
\begin{tablenotes}Each statement was checked symbolically (S: reduced to zero by computer algebra and evaluated at random
rational points), numerically (N: closed forms rederived without using them, by brute-force optimization, numerically
traced isoquants, finite-difference Jacobians, or direct Fisher-information computation), or by Monte Carlo (MC: at
least $4\times10^5$ draws). The register contains 61 checks, all of which pass. ``Corrected'' marks a statement whose
earlier formulation, in our working drafts or in common readings of the literature, was imprecise or false; the table
states the corrected version, which is the one used in the paper. Proofs are in Online Appendix~\ref{app:proofs}. The
information, selection, and transmission statements assume Gaussian noise where closed forms are given; the
inequalities are distribution-free.\end{tablenotes}
\begin{tablenotes}[Source]Authors' derivations and verification code.\end{tablenotes}
\end{table}


\clearpage
